\documentclass[11pt,reqno]{amsart}

% ================================================================
%  Packages
% ================================================================
\usepackage{amsmath,amssymb,amsthm}
\usepackage{mathrsfs}
\usepackage{enumerate}
\usepackage[shortlabels]{enumitem}
\usepackage{booktabs}
\usepackage{array}
\usepackage{microtype}
\usepackage[dvipsnames]{xcolor}
\usepackage[colorlinks=true,linkcolor=blue!60!black,citecolor=green!40!black,urlcolor=blue!60!black]{hyperref}
\usepackage[all]{xy}

% ================================================================
%  Theorem environments
% ================================================================
\newtheorem{theorem}{Theorem}[section]
\newtheorem{proposition}[theorem]{Proposition}
\newtheorem{lemma}[theorem]{Lemma}
\newtheorem{corollary}[theorem]{Corollary}
\newtheorem{conjecture}[theorem]{Conjecture}
\theoremstyle{definition}
\newtheorem{definition}[theorem]{Definition}
\newtheorem{construction}[theorem]{Construction}
\newtheorem{example}[theorem]{Example}
\theoremstyle{remark}
\newtheorem{remark}[theorem]{Remark}

% ================================================================
%  Macros
% ================================================================
\newcommand{\cA}{\mathcal{A}}
\newcommand{\cM}{\mathcal{M}}
\newcommand{\barB}{\bar{B}}
\newcommand{\barC}{\overline{C}}
\newcommand{\Defcyc}{\mathrm{Def}_{\mathrm{cyc}}}
\newcommand{\Ran}{\mathrm{Ran}}
\newcommand{\MC}{\mathrm{MC}}
\newcommand{\Sym}{\mathrm{Sym}}
\newcommand{\Hom}{\mathrm{Hom}}
\newcommand{\Ext}{\mathrm{Ext}}
\newcommand{\Res}{\mathrm{Res}}
\newcommand{\Aut}{\mathrm{Aut}}
\newcommand{\FM}{\overline{C}}
\newcommand{\fg}{\mathfrak{g}}
\newcommand{\fh}{\mathfrak{h}}
\newcommand{\gmod}{\mathfrak{g}^{\mathrm{mod}}_\cA}
\newcommand{\Sh}{\mathrm{Sh}}
\newcommand{\OPE}{\mathrm{OPE}}
\newcommand{\Vir}{\mathrm{Vir}}
\newcommand{\Ahat}{\widehat{A}}
\newcommand{\cW}{\mathcal{W}}

\DeclareMathOperator{\gr}{gr}
\DeclareMathOperator{\rk}{rk}
\DeclareMathOperator{\depth}{depth}

\numberwithin{equation}{section}

% ================================================================
\begin{document}

\title[Modular Koszul Duality and the Shadow Tower]
{Modular Koszul Duality and the Shadow Tower:\\
Algebraicity, Depth Classification,\\
and the Gravitational Coproduct}

\author{Raeez Lorgat}
\address{Department of Mathematics, Harvard University, Cambridge, MA 02138}
\email{raeez@math.harvard.edu}

\date{\today}

\begin{abstract}
Let $\cA$ be a chirally Koszul vertex algebra on a smooth projective
curve~$X$, with modular characteristic $\kappa = \kappa(\cA)$, cubic
shadow $\alpha$, and quartic shadow~$S_4$.  The bar construction
$\barB^{\mathrm{ch}}(\cA)$ defines a Maurer--Cartan element
$\Theta_\cA$ in the modular convolution dg Lie algebra, whose
finite-order projections $\Theta_\cA^{\le r}$ form the
\emph{shadow Postnikov tower}.

We prove six results.
First, on each one-dimensional primary line~$L$ in the
cyclic deformation complex, the weighted shadow generating function
$H(t) = \sum_{r \ge 2} r\, S_r\, t^r$ satisfies the algebraic
relation $H(t)^2 = t^4 Q(t)$, where
$Q(t) = 4\kappa^2 + 12\kappa\alpha\, t
+ (9\alpha^2 + 16\kappa S_4)\, t^2$
is the \emph{shadow metric}: the tower is algebraic of degree~$2$
along~$L$, and every coefficient $S_r$ for $r \ge 5$ is determined
by the three seed invariants $(\kappa, \alpha, S_4)$ of~$L$.
Second, the \emph{critical discriminant} $\Delta = 8\kappa S_4$
partitions the standard Lie-theoretic landscape into four depth
classes: Gaussian ($r_{\max} = 2$), Lie ($r_{\max} = 3$), contact
($r_{\max} = 4$), and mixed ($r_{\max} = \infty$), with $\Delta = 0$
on the primary line forcing finite depth and $\Delta \ne 0$
forcing infinite depth.
Third, we compute the genus-$2$ free energy
$F_2 = \kappa \cdot 7/5760$ for single-generator algebras
by an explicit stable-graph sum, carry out the first
multi-channel computation for~$\mathcal{W}_3$ (exhibiting
a cross-channel correction whose cancellation is equivalent
to multi-generator universality at genus~$2$), and derive
the planted-forest correction
$\delta_{\mathrm{pf}}^{(2,0)} = S_3(10 S_3 - \kappa)/48$.
Fourth, we prove that the gravitational coproduct, obtained by
homological perturbation transfer through BRST reduction from
$\widehat{\mathfrak{sl}}_N$ to $\cW_N$, is strictly primitive
at all arities for every principal Drinfeld--Sokolov reduction:
gravity does not produce particles, only braids.
Fifth, for class~$\mathbf{M}$ algebras at rational central
charge, the shadow metric~$Q_L$ defines an Epstein zeta function
$\varepsilon_{Q_L}(s) = \sum'_{(m,n)} Q_L(m,n)^{-s}$
with a functional equation $s \leftrightarrow 1-s$, and the
discriminant of~$Q_L$ determines an imaginary quadratic
\emph{shadow field}
$K_L = \mathbb{Q}(\sqrt{\operatorname{disc}(Q_L)})$
through which the Epstein zeta factors via a
Dirichlet $L$-function: the chain
$\cA \to Q_L \to K_L \to L(s, \chi_d)$ produces
arithmetic from OPE data.
Sixth, we verify the genus-$2$ prediction $F_2(V_{E_8}) = 7/720$
against the Siegel--Weil formula and direct lattice enumeration:
the bar complex knows that each of the~$240$ $E_8$ roots has
exactly~$126$ orthogonal partners.
\end{abstract}

\subjclass[2020]{17B69, 14H10, 18M70, 81T40}

\keywords{Chiral algebras, bar construction, Koszul duality,
moduli of curves, vertex algebras, Maurer--Cartan,
W-algebras, gravitational coproduct, Epstein zeta function,
shadow field}

\maketitle

% ================================================================
%  1. INTRODUCTION
% ================================================================

\section{Introduction}\label{sec:intro}

\subsection{The question}

Classical Koszul duality is a phenomenon of graded algebras
over a point.  A quadratic algebra~$A$ determines a dual
coalgebra~$A^!$; the bar construction $B(A)$ mediates between
them; and the Koszul property (bar cohomology concentrated in
bar degree~$1$) is equivalent to the existence of a linear
resolution.  The theory is the subject of Priddy~\cite{Priddy70},
Beilinson--Ginzburg--Soergel~\cite{BGS96}, and
Loday--Vallette~\cite{LV12}.

On an algebraic curve~$X$, the generators of the algebra become
sections of a $\mathcal{D}_X$-module, the quadratic relations
become operator product expansions with meromorphic singularities,
and the bar differential becomes an integral transform whose
kernel is the logarithmic $1$-form
$\eta_{ij} = d\log(z_i - z_j)$ on the Fulton--MacPherson
compactification $\FM_n(X)$ of the configuration space.  The
resulting object, the \emph{chiral bar complex}
$\barB^{\mathrm{ch}}(\cA)$, is a factorization coalgebra on the
Ran space $\Ran(X)$ \cite{BD04, FG12, CG17}.

At genus~$0$, the classical theory embeds faithfully.  At
genus~$g \ge 1$, the topology of the curve forces genuinely new
structure: the bar differential acquires curvature
$d_{\mathrm{fib}}^2 = \kappa(\cA) \cdot \omega_g$, where
$\omega_g = c_1(\mathbb{E})$ is the Hodge class on
$\overline{\cM}_g$ and the scalar
$\kappa(\cA) \in \mathbb{Q}(c)$ is an intrinsic invariant of
the algebra called the \emph{modular characteristic}.  This
curvature has no analogue in the classical setting; its geometric
origin is the nontriviality of the Hodge bundle over
$\overline{\cM}_g$.

The present paper addresses a single question: \emph{what
invariants does the chiral bar construction produce, beyond the
modular characteristic?}

\subsection{The answer: the shadow obstruction tower}

The bar complex on a curve of genus~$g$ with $n$ marked points
defines a pronilpotent dg Lie algebra $\gmod$ (the modular
convolution algebra), whose differential encodes graph sums over
stable curves.  The boundary-of-boundary relation
$\partial^2 = 0$ on $\overline{\cM}_{g,n}$ implies $D^2 = 0$ on
$\gmod$, and the element $\Theta_\cA := D_\cA - d_0$ is
therefore a Maurer--Cartan element:
$D \Theta_\cA + \frac{1}{2}[\Theta_\cA, \Theta_\cA] = 0$.

The MC element $\Theta_\cA$ encodes the full genus expansion.
Its finite-order projections $\Theta_\cA^{\le r}$, obtained by
truncating the arity filtration at level~$r$, form the
\emph{shadow Postnikov tower}.  At arity~$2$, the tower
produces the modular characteristic~$\kappa$.  At arity~$3$,
a cubic shadow~$\alpha$ (the structure constant of the transferred
$L_\infty$ bracket $\ell_3$).  At arity~$4$, a quartic
shadow~$S_4$ (the first genuinely nonlinear invariant, measuring
the contact term in the four-point collision residue).

The obstruction to truncating $\Theta_\cA$ at finite arity is
measured by classes
$o_{r+1}(\cA) \in H^2(\cA^{\mathrm{sh}}_{r+1,0})$ in the
shadow algebra $\cA^{\mathrm{sh}} := H_\bullet(\Defcyc^{\mathrm{mod}}(\cA))$.
When $o_{r+1} = 0$, the tower stops; when $o_{r+1} \ne 0$, the
next shadow coefficient $S_{r+1}$ is forced.  The question is
how far this process continues and what controls it.

\subsection{Five movements}

The paper develops in five movements.  First, the universal
MC element $\Theta_\cA$ and its shadow obstruction tower as the primary
object.  Second, the algebraicity theorem: the tower is a
square root, $H = t^2\sqrt{Q_L}$, controlled by a quadratic
form in the arity parameter.  Third, the four-class partition
G/L/C/M determined by the critical discriminant $\Delta$.
Fourth, the explicit computations: genus-$2$ graph sums,
the $E_8$ decisive test, gravitational coproduct primitivity.
Fifth, the arithmetic frontier: the Epstein zeta function of
the shadow metric, the shadow field, and the descent fan from
$\Theta_\cA$ to three mathematical worlds.

\subsection{The shadow metric and algebraicity}

The first main result (Theorem~\ref{thm:main}, parts (i)--(ii))
is that the shadow obstruction tower is controlled by a single quadratic
polynomial.

Restrict to a one-dimensional primary slice~$L$ of the
cyclic deformation complex (the space spanned by a single primary
field~$\eta$ of definite conformal weight).  On~$L$, the
MC recursion reduces to a scalar recursion for the shadow
coefficients $S_r$, and the weighted generating function
$H(t) = \sum_{r \ge 2} r\,S_r\,t^r$ satisfies the algebraic
relation
\begin{equation}\label{eq:intro-algebraic}
  H(t)^2 \;=\; t^4\,Q(t),
\end{equation}
where the \emph{shadow metric}
\begin{equation}\label{eq:intro-shadow-metric}
  Q(t)
  \;=\;
  4\kappa^2 + 12\kappa\alpha\,t
  + (9\alpha^2 + 16\kappa S_4)\,t^2
\end{equation}
is a quadratic polynomial whose coefficients are determined by the
three seed invariants $(\kappa, \alpha, S_4)$ of the algebra.
The generating function is therefore algebraic of degree~$2$:
$H(t) = t^2 \sqrt{Q(t)}$.  Every shadow coefficient $S_r$
for $r \ge 5$ is a rational function of $(\kappa, \alpha, S_4)$,
computable from the binomial expansion of~$\sqrt{Q}$.

The proof converts the Maurer--Cartan recursion into a convolution
identity for the Taylor coefficients of $F(t) = H(t)/t^2$,
showing $F^2 = Q$.  The orders $m = 0, 1, 2$ of this identity
reproduce the definitions of $\kappa$, $\alpha$, $S_4$; the
orders $m \ge 3$ are exactly the recursion.

This is the \emph{intrinsic quartic principle}: an infinite
sequence of invariants, one for each arity $r \ge 2$, is
determined by three numbers.  The situation is analogous to the
role of the Weierstrass equation for an elliptic curve: two
parameters $(g_2, g_3)$ determine all Laurent coefficients
of~$\wp$, because $\wp$ satisfies an algebraic equation of
degree~$2$ in $\wp'$.  Here, three parameters
$(\kappa, \alpha, S_4)$ determine all shadow coefficients
$S_r$, because $H$ satisfies an algebraic equation of
degree~$2$ in~$H$.

\subsection{The four-class classification}

Write
$Q(t) = (2\kappa + 3\alpha t)^2 + 2\Delta\,t^2$
for the \emph{Gaussian decomposition}, where
$\Delta := 8\kappa S_4$ is the \emph{critical discriminant}.
The Gaussian envelope $(2\kappa + 3\alpha t)^2$ is the part
of~$Q$ visible to the quadratic and cubic shadows alone;
$\Delta$ measures the irreducible quartic contribution.

The critical discriminant controls a fundamental dichotomy.

When $\Delta = 0$, the shadow metric is a perfect square:
$Q = (2\kappa + 3\alpha t)^2$, the generating function
$H(t) = t^2(2\kappa + 3\alpha t)$ is polynomial, and the
tower terminates.  If $\alpha = 0$, only $S_2 = \kappa$ survives
(depth~$2$); if $\alpha \ne 0$, then $S_2$ and $S_3$ survive
(depth~$3$).

When $\Delta \ne 0$, the shadow metric is irreducible:
$\sqrt{Q}$ is irrational, every Taylor coefficient of $\sqrt{Q}$
beyond degree~$1$ carries a factor of~$\Delta$, and $S_r \ne 0$
for all $r \ge 4$ (generically).  The tower is infinite.

This yields the second main result
(Theorem~\ref{thm:main}, part~(iii)): every chirally Koszul
algebra whose primary-line shadow data has been computed falls
into one of four \emph{shadow depth classes}.  The partition
is exhaustive for the standard Lie-theoretic landscape
(Table~\ref{tab:master-shadow}); whether every chirally Koszul
algebra belongs to one of these four classes is a structural
observation, not a classification theorem.

\begin{center}
\small
\renewcommand{\arraystretch}{1.3}
\begin{tabular}{@{}lccll@{}}
\toprule
Class & $\Delta$ & $r_{\max}$ & Prototype & Shadow data \\
\midrule
$\mathbf{G}$ (Gaussian) & $0$ & $2$ & Heisenberg & $\kappa = k$ \\
$\mathbf{L}$ (Lie) & $0$ & $3$ & Affine Kac--Moody &
$\kappa,\,\alpha \ne 0$ \\
$\mathbf{C}$ (Contact) & $0^*$ & $4$ & $\beta\gamma$ system &
$\kappa = 6\lambda^2{-}6\lambda{+}1$ \\
$\mathbf{M}$ (Mixed) & $\ne 0$ & $\infty$ & Virasoro, $\mathcal{W}_N$ &
$\kappa,\,\alpha,\,S_4$ \\
\bottomrule
\multicolumn{5}{@{}l}{\footnotesize ${}^*$Class $\mathbf{C}$:
$\Delta = 0$ on the primary line; the quartic contact invariant
lives on a charged stratum.}
\end{tabular}
\end{center}

\noindent
Class~$\mathbf{G}$ consists of algebras whose bar
complex is formal in the classical sense: the shadow obstruction tower
terminates at the leading scalar, and all higher-genus free
energies $F_g$ are determined by~$\kappa$ alone.
Class~$\mathbf{L}$ adds exactly one cubic correction (the
structure constant of the Lie bracket on generators), after which
the tower terminates by the Jacobi identity.
Class~$\mathbf{C}$ is the first nonlinear class: the quartic
contact invariant~$S_4$ is nonzero but lives on a separate
stratum whose self-bracket exits the complex, so the tower
terminates after one further step.
Class~$\mathbf{M}$ is the generic class: the shadow obstruction tower is
infinite, the free energies $F_g$ receive corrections at every
arity, and the shadow coefficients $S_r$ decay (or grow) at a
computable rate governed by the \emph{shadow growth rate}
\begin{equation}\label{eq:intro-growth-rate}
  \rho(\cA)
  \;=\;
  \sqrt{\frac{9\alpha^2 + 16\kappa S_4}{4\kappa^2}}.
\end{equation}

The growth rate is the reciprocal of the branch-point modulus
of the shadow metric.  When $\rho < 1$, the shadow obstruction tower converges
absolutely; when $\rho > 1$, it diverges but the algebraic
function $\sqrt{Q}$ extends analytically past the convergence
radius.  For the Virasoro algebra at central charge~$c$:
\begin{equation}\label{eq:intro-virasoro-rho}
  \rho(\Vir_c)
  \;=\;
  \frac{1}{c}\sqrt{\frac{180c + 872}{5c + 22}},
\end{equation}
with a critical central charge $c_\star \approx 6.125$ (the
unique positive root of $5c^3 + 22c^2 - 180c - 872 = 0$)
separating the convergent and divergent regimes.  At the
self-dual point $c = 13$ (where $\Vir_c^! = \Vir_{26-c}$ satisfies
$\Vir_c = \Vir_c^!$): $\rho \approx 0.467$.

Classical Koszul duality over a point is the genus-$0$,
arity-$2$, $\Delta = 0$ stratum of this picture: the formal
Gaussian case where the tower terminates immediately.  Everything
beyond class~$\mathbf{G}$ is new.

\subsection{The shadow connection and Koszul monodromy}

The shadow metric equips each primary line with a logarithmic
connection
\begin{equation}\label{eq:intro-connection}
  \nabla^{\mathrm{sh}}
  \;=\;
  d \;-\; \frac{Q'(t)}{2\,Q(t)}\,dt,
\end{equation}
with regular singularities at the zeros of~$Q$ and residue
$\frac{1}{2}$.  The monodromy around each branch point is~$-1$.
This is the \emph{Koszul sign}: the square root $\sqrt{Q}$
changes sign around a simple zero, and the shadow obstruction tower inherits
this sign change through the algebraic relation
$H = t^2\sqrt{Q}$.

Under Koszul duality $\cA \mapsto \cA^!$, the shadow metric
transforms by the Verdier involution.  For the Virasoro family
$\Vir_c^! = \Vir_{26-c}$, the discriminants satisfy
\begin{equation}\label{eq:intro-disc-complementarity}
  \Delta(\Vir_c) + \Delta(\Vir_{26-c})
  \;=\;
  \frac{6960}{(5c + 22)(152 - 5c)},
\end{equation}
a Virasoro-specific identity whose numerator
$6960 = 40 \cdot 174$ is a product of the Virasoro quartic
normalization and the sum of the Lee--Yang denominators.
The Koszul dual singularities $c = -22/5$ and $c = 152/5$
sum to~$26$: they are Koszul partners.  For other families
(affine Kac--Moody, $\cW_N$ with $N \ge 3$), the
discriminant sum has a different numerator.

\subsection{Genus-$2$ free energies}

The shadow obstruction tower produces invariants of moduli spaces by
evaluating graph sums over stable curves.  The genus-$g$
\emph{free energy}
$F_g(\cA) = \kappa(\cA) \cdot \lambda_g^{\mathrm{FP}}$
is computed for each chirally Koszul algebra by the
\emph{Faber--Pandharipande integral}
$\lambda_g^{\mathrm{FP}} = \int_{\overline{\cM}_g} \lambda_g$,
where $\lambda_g = c_g(\mathbb{E})$ is the top Chern class
of the Hodge bundle.

At genus~$2$: $\lambda_2^{\mathrm{FP}} = 7/5760$, and
$F_2(\cA) = \kappa \cdot 7/5760$ for every single-generator
chirally Koszul algebra.  The verification requires
an explicit sum over the six stable graphs of
genus~$2$ with no marked points (the smooth genus-$2$
curve, figure-eight, banana, dumbbell, theta, and
lollipop), weighted by automorphism factors and vertex
contributions from the shadow obstruction tower.

For the $\mathcal{W}_3$ algebra, which has two generators
$T$ (weight~$2$) and $W$ (weight~$3$) producing two primary
lines, the genus-$2$ computation is a multi-channel graph sum:
each edge of the stable graph carries a channel label
($T$ or~$W$), each vertex contributes a shadow coefficient on
the corresponding primary line, and the propagator
$P_T = 2/c$, $P_W = 3/c$ depends on the channel.  The
per-channel diagonal sum yields
$F_2^{\mathrm{scal}}(\mathcal{W}_3) = \kappa \cdot \lambda_2^{\mathrm{FP}}
= 7c/6912$, but the full multi-channel graph sum produces an
additional cross-channel correction
$\delta F_2 = (c+120)/(16c)$.  Whether this correction is
cancelled by CohFT $R$-matrix terms (restoring the universality
identity $F_g = \kappa \cdot \lambda_g^{\mathrm{FP}}$ at
genus~$2$ for multi-generator algebras) is the content of the
multi-generator universality problem, which remains open.

The Maurer--Cartan equation on
$\overline{\cM}_{2,0}$ also produces a
\emph{planted-forest correction}
\begin{equation}\label{eq:intro-planted-forest}
  \delta_{\mathrm{pf}}^{(2,0)}
  \;=\;
  \frac{S_3(10\,S_3 - \kappa)}{48},
\end{equation}
measuring the contribution of codimension-$2$ boundary strata
to the genus-$2$ amplitude.  This correction vanishes for
class~$\mathbf{G}$ (where $S_3 = 0$).  For
classes~$\mathbf{C}$ and~$\mathbf{M}$, it is generically nonzero.
For the Virasoro algebra:
$\delta_{\mathrm{pf}}^{(2,0)} = -(c - 40)/48$.

\subsection{The gravitational coproduct}

The fourth result concerns the transferred algebraic structure
under BRST reduction.  The Drinfeld--Sokolov (DS) reduction
functor takes affine Kac--Moody data to $\cW$-algebra data via a
BRST complex with ghosts.  The homological perturbation lemma
(HPL) transfers the $A_\infty$ products and the coproduct from
the affine dg-shifted Yangian to a transferred Yangian on the
$\cW$-algebra sector.

We prove the result for all principal DS reductions
$\widehat{\mathfrak{sl}}_N \to \cW_N$.
The ghost-number obstruction (the homotopy~$h$ in the SDR has
ghost degree~$-1$, and the projection $p \otimes p$ kills all
ghost-number-$(-1)$ elements) forces the transferred coproduct to
be \emph{strictly primitive at all arities}.  The mechanism is
universal: the ghost-number grading of the BRST complex for any
principal DS reduction $\widehat{\fg} \to \cW(\fg)$ has the same
structure.  For $N = 2$ (the Virasoro case):
\[
  \Delta_z^{\Vir}(T)
  \;=\;
  \tau_z(T) \otimes 1 + 1 \otimes T,
  \qquad
  \Delta_{z,n}^{\Vir} = 0
  \quad\text{for } n \ge 2.
\]
All higher-arity coproduct corrections, which would represent
graviton splitting into constituent gravitons, are killed by the
ghost-number mechanism.  The only nontrivial gravitational
structure is the $r$-matrix twist
$r^{\Vir}(z) = (c/2)/z^3 + 2T/z$ on the target product,
encoding the braiding of line operators.

The physical content: in Chern--Simons gauge theory, gluons
split (the coproduct is nontrivial, reflecting the
Clebsch--Gordan decomposition of tensor products).  In
$3$d gravity, gravitons do not split.  The gravitational
coproduct is primitive; all interactions are encoded in
the twist, not in particle production.

\subsection{The Epstein zeta function of the shadow metric}

For class~$\mathbf{M}$ algebras at rational central charge,
the shadow metric~$Q_L$ is a positive definite binary quadratic
form on~$\mathbb{Z}^2$ (since
$\operatorname{disc}(Q_L) = -32\kappa^2\Delta < 0$).  Its
Epstein zeta function
\[
  \varepsilon_{Q_L}(s)
  \;:=\;
  \sideset{}{'}\sum_{(m,n) \in \mathbb{Z}^2}
  Q_L(m,n)^{-s}
\]
converges for $\operatorname{Re}(s) > 1$, extends
meromorphically, and satisfies the functional equation
$\Lambda_{Q_L}(s) = \Lambda_{Q_L}(1 - s)$ after the standard
gamma completion~\cite{Epstein1903, Terras73}.  The
discriminant of~$Q_L$ determines an imaginary quadratic
extension $K_L = \mathbb{Q}(\sqrt{D_0})$, the \emph{shadow
field}.  When the class number $h(D_0) = 1$, the Epstein
zeta factors through the Dedekind zeta of~$K_L$, which in
turn splits as $\zeta_{K_L}(s) = \zeta(s) \cdot L(s, \chi_d)$.
The chain
\begin{equation}\label{eq:intro-shadow-field-chain}
  \cA
  \;\xrightarrow{\;\text{shadow obstruction tower}\;}
  Q_L
  \;\xrightarrow{\;\operatorname{disc}\;}
  D_0
  \;\xrightarrow{\;\mathbb{Q}(\sqrt{\cdot})\;}
  K_L
  \;\xrightarrow{\;\zeta_K\;}
  L(s, \chi_d)
\end{equation}
produces a Dirichlet $L$-function from OPE data.  For the
Ising model ($c = 1/2$): the shadow field is
$\mathbb{Q}(\sqrt{-10})$.  For class~$\mathbf{G}$ and
class~$\mathbf{L}$ algebras, $\Delta = 0$ forces the
discriminant to vanish and the construction degenerates.
This is the shadow-theoretic meaning of the G/L/C/M
partition: classes~$\mathbf{G}$ and~$\mathbf{L}$ have
trivial arithmetic content, while class~$\mathbf{M}$ produces
genuine number-theoretic invariants.

\subsection{The $E_8$ decisive test}

The bar complex makes a prediction for every lattice vertex
algebra at every genus.  At genus~$2$ for the $E_8$ lattice:
$F_2(V_{E_8}) = 8 \cdot 7/5760 = 7/720$.  This prediction
is verified by three independent methods:
(a)~the bar-complex formula $F_2 = \kappa \cdot \lambda_2^{\mathrm{FP}}$;
(b)~the Siegel--Weil formula with Cohen--Katsurada coefficients
($\Theta_{E_8}^{(2)} = E_4^{(2)}$, giving
$a(I_2) = 240 \times 126 = 30240$);
(c)~direct lattice enumeration.
The three-way agreement is a nontrivial consistency check:
the bar complex knows, from purely algebraic data, that each
of the $240$ $E_8$ roots has exactly~$126$ orthogonal partners.

\subsection{The descent fan}

The MC element $\Theta_\cA$ projects onto three mathematical
worlds.  These projections share a common source but do not
compose into a single functor: the ``descent chain'' is a
\emph{fan}, not a chain.

\begin{enumerate}[(i)]
\item \emph{Categorical projection.}
  $\cA \mapsto \barB^{\mathrm{ch}}(\cA) \mapsto \cA^!$.
  Target: factorization coalgebras and their Verdier duals.
  Symmetry: the Koszul involution $\cA \leftrightarrow \cA^!$.

\item \emph{Spectral projection.}
  $\cA \mapsto \{S_r\} \mapsto \sqrt{Q_L} \mapsto
  \varepsilon_{Q_L}(s)$.
  Target: the shadow obstruction tower and the Epstein zeta.
  Symmetry: the functional equation $s \leftrightarrow 1 - s$.

\item \emph{Modular projection.}
  $\cA \mapsto Z_g(\cA) \mapsto S_\cA(v)$.
  Target: the genus expansion and the sewing partition function.
  Symmetry: the Galois involution on the $L$-function content.
\end{enumerate}
The Koszul involution, the functional equation involution, and
the Galois involution are three different involutions on three
different spaces.  No functor intertwines all three.

\subsection{What is proved, what is open}

The following is a candid assessment.

\emph{Proved unconditionally}: the algebraicity theorem
(Theorem~\ref{thm:main}(i)--(ii)); the four-class
classification; the genus-$g$ formula
$F_g = \kappa \cdot \lambda_g^{\mathrm{FP}}$ for all
single-generator chirally Koszul algebras at all genera;
the genus-$1$ universality $F_1 = \kappa/24$ for all
families including multi-generator; the gravitational
coproduct primitivity for all principal DS reductions;
the Epstein zeta functional equation for class~$\mathbf{M}$
at rational~$c$.

\emph{Open}: whether $F_g = \kappa \cdot \lambda_g^{\mathrm{FP}}$
holds at genus~$g \ge 2$ for multi-generator algebras
(the multi-generator universality problem).  The
genus-$2$ $\cW_3$ computation
(\S\ref{sec:genus2-w3}) exhibits the cross-channel
correction that any proof must cancel.  Whether
the Koszul--Epstein zeta has zeros on a single critical
line is completely open; it requires understanding the
arithmetic of the shadow field at a level far beyond
the present methods.

\subsection{Relation to existing work}

The bar construction for chiral algebras was introduced by
Beilinson and Drinfeld~\cite{BD04} and developed in the
factorization algebra framework by Francis and Gaitsgory
\cite{FG12} and Costello and Gwilliam~\cite{CG17}.
The modular characteristic~$\kappa$ appears implicitly in the
work of Feigin and Fuchs on Hodge bundles of conformal blocks,
and the formula $\kappa(\Vir_c) = c/2$ can be extracted from
Beilinson and Schechtman~\cite{BS88}.
The invariants of moduli spaces produced by vertex algebras via
conformal blocks are the subject of Faber and
Pandharipande~\cite{FP00}, Mumford~\cite{Mumford83}, and the
extensive literature on cohomological field theories following
Kontsevich and Manin~\cite{KM94}.  The Epstein zeta function
of a binary quadratic form is due to Epstein~\cite{Epstein1903};
the functional equation and Bessel-series expansion are due
to Terras~\cite{Terras73}.  The connection between conformal
field theory and the scalar modular bootstrap is studied
by Benjamin and Chang~\cite{BenjaminChang22}.

The quartic shadow $S_4$, the critical discriminant~$\Delta$,
the four-class partition, the algebraicity theorem
(Theorem~\ref{thm:main}), the shadow field, and the
Koszul--Epstein construction appear to be new.  The connection
between shadow depth and $L_\infty$ non-formality, while
consonant with the general philosophy of Kontsevich~\cite{Kontsevich03},
has not been established previously in the chiral setting.

The Drinfeld--Sokolov reduction and its homological algebra are
classical \cite{KRW03, FF90, DS85};
the gravitational coproduct primitivity result
(Theorem~\ref{thm:grav-primitivity-standalone}) and its connection to the
shadow obstruction tower appear to be new.

\subsection{Structure of the paper}

Section~\ref{sec:background} reviews chiral algebras, the bar
construction, and the modular characteristic, establishing
notation and the three seed invariants
$(\kappa, \alpha, S_4)$.
Section~\ref{sec:convolution} defines the modular convolution
algebra and the Maurer--Cartan element~$\Theta_\cA$.
Section~\ref{sec:shadow-metric} introduces the shadow
metric and the recursion on primary lines.
Section~\ref{sec:algebraicity} proves the Riccati algebraicity
theorem, derives the shadow connection, the depth
classification, and the growth rate.
Section~\ref{sec:four-classes} presents the four shadow
classes with explicit examples and the master shadow table.
Section~\ref{sec:shadow-tables} records the complete shadow
tables for the standard landscape, including the
complementarity data and the Virasoro tower through arity~$10$.
Section~\ref{sec:genus2-w3} carries out the genus-$2$ graph sum
for~$\cW_3$, including the multi-channel cross-channel correction
and the planted-forest formula.
Section~\ref{sec:gravitational-coproduct} proves the gravitational
coproduct primitivity for all principal DS reductions.
Section~\ref{sec:e8-test} verifies the $E_8$ decisive test.
Section~\ref{sec:epstein} develops the Epstein zeta function
of the shadow metric and the shadow field construction.
Section~\ref{sec:outlook} discusses the descent fan, the
Koszul--Epstein frontier, and open problems.


% ================================================================
%  Body sections
% ================================================================

\section{Chiral Koszul pairs and the bar construction}
\label{sec:background}

Let $X$ be a smooth projective curve of genus~$g$ over an
algebraically closed field~$k$ of characteristic zero.
A \emph{chiral algebra} on~$X$ is a $\mathcal{D}_X$-module
$\cA$ equipped with a chiral bracket
$\mu \colon j_* j^* (\cA \boxtimes \cA)
\to \Delta_! \cA$, where $j \colon X^2 \setminus \Delta
\hookrightarrow X^2$ is the complement of the diagonal and
$\Delta_!$ is the $!$-pushforward~\cite{BD04}.  On the formal
disc, the chiral bracket encodes the operator product expansion
$\phi^a(z)\,\phi^b(w) \sim \sum_{n \ge 0}
(\phi^a{}_{(n)}\phi^b)(w)/(z-w)^{n+1}$.

The \emph{chiral bar construction}
$\barB^{\mathrm{ch}}(\cA) = \bigoplus_{n \ge 1}
\cA^{\boxtimes n}[-1]^{\otimes n}$ is a factorization coalgebra
on the Ran space $\Ran(X)$, with differential
$d = \sum_{i < j} d_{ij}$ built from the OPE residues
weighted by the kernel $\eta_{ij} = d\log(z_i - z_j)$ on the
Fulton--MacPherson compactification~\cite{FG12, CG17}.
The bar complex is \emph{chirally Koszul} if
$H^*(\barB^{\mathrm{ch}}(\cA))$ is concentrated in bar
degree~$1$.

The \emph{modular characteristic} of~$\cA$ is the scalar
$\kappa(\cA) \in \mathbb{Q}(c)$ defined by the genus-$1$
identity $F_1(\cA) = \kappa(\cA)/24$.  It controls the
bar curvature: $d^2|_{\text{fib}} = \kappa \cdot \omega_g$,
where $\omega_g = c_1(\mathbb{E})$ is the Hodge class.
Explicit formulas: $\kappa(\mathcal{H}_k) = k$,
$\kappa(V_k(\fg)) = \dim(\fg)(k + h^\vee)/(2h^\vee)$,
$\kappa(\Vir_c) = c/2$, $\kappa(\cW_N) = c(H_N - 1)$.

The \emph{cubic shadow} $\alpha = S_3$ measures the transferred
$L_\infty$ bracket~$\ell_3$ on bar cohomology: it is the
structure constant of the three-point collision residue.  The
\emph{quartic shadow} $S_4$ is the contact term of the
four-point collision, computed from the quartic OPE data.
These three numbers $(\kappa, \alpha, S_4)$ are the seed
invariants of the paper.

\section{The modular convolution algebra}
\label{sec:convolution}

The bar complex on a curve of genus~$g$ with~$n$ marked points
defines a pronilpotent dg Lie algebra
$\gmod$ (the modular convolution algebra), whose differential
encodes graph sums over stable curves.  The construction proceeds
in three steps.

\emph{Step 1: the coefficient algebra.}
Let $G_{\mathrm{st}}$ denote the species of stable graphs
$\Gamma = (V, E, g \colon V \to \mathbb{Z}_{\ge 0})$ with
genus function~$g$ and the stability condition
$2g_v + \mathrm{val}_v \ge 3$ at each vertex.  The stable-graph
coefficient algebra is
$\mathfrak{G}_{\mathrm{st}} = \bigoplus_\Gamma
\mathbb{Q} \cdot [\Gamma] / \mathrm{Aut}(\Gamma)$,
graded by $h^1(\Gamma) + \sum_v g_v$.

\emph{Step 2: convolution.}
The modular convolution dg Lie algebra is
$\gmod = \Hom(\mathfrak{G}_{\mathrm{st}}, \Defcyc(\cA))$,
with bracket from graph composition (inserting a graph at a
vertex of another) and differential from the boundary-of-boundary
relation $\partial^2 = 0$ on $\overline{\cM}_{g,n}$.

\emph{Step 3: the MC element.}
The boundary relation $\partial^2 = 0$ implies $D^2 = 0$ on
$\gmod$.  The element $\Theta_\cA := D_\cA - d_0$ is
therefore a Maurer--Cartan element:
$D\Theta_\cA + \tfrac{1}{2}[\Theta_\cA, \Theta_\cA] = 0$.
Its finite-order projections $\Theta_\cA^{\le r}$ form the
shadow Postnikov tower.

\begin{theorem}[Main theorem]\label{thm:main}
Let\/ $\cA$ be a chirally Koszul vertex algebra on a smooth
projective curve with modular characteristic $\kappa$,
cubic shadow~$\alpha$, and quartic shadow~$S_4$.
\begin{enumerate}[label=\textup{(\roman*)}]
\item \textbf{Algebraicity}
  \textup{(}Theorem~\textup{\ref{thm:riccati-algebraicity}}\textup{)}.
  $H(t)^2 = t^4 Q_L(t)$, where
  $Q_L = 4\kappa^2 + 12\kappa\alpha\,t
  + (9\alpha^2 + 16\kappa S_4)\,t^2$.
\item \textbf{Classification}
  \textup{(}Theorem~\textup{\ref{thm:single-line-dichotomy}}\textup{)}.
  $r_{\max} \in \{2, 3, 4, \infty\}$.
\item \textbf{Gravitational primitivity}
  \textup{(}Theorem~\textup{\ref{thm:grav-primitivity-standalone}}\textup{)}.
  $\Delta_{z,k}^{\cW(\fg)} = 0$ for all $k \ge 2$.
\item \textbf{Epstein zeta}
  \textup{(}Theorem~\textup{\ref{thm:shadow-epstein}}\textup{)}.
  $\Lambda_{Q_L}(s) = \Lambda_{Q_L}(1-s)$.
\item \textbf{$E_8$ test}
  \textup{(}Theorem~\textup{\ref{thm:e8-three-way}}\textup{)}.
  $F_2(V_{E_8}) = 7/720$.
\end{enumerate}
\end{theorem}

% riccati.tex — §3–4 of the standalone paper
% The shadow metric, the Riccati algebraicity theorem, and the depth classification.

% ================================================================
% §3. The Shadow Metric and the Riccati Equation
% ================================================================

\section{The shadow metric and the Riccati equation}
\label{sec:shadow-metric}

\subsection{The shadow obstruction tower}

Let $\mathcal{A}$ be a modular Koszul chiral algebra with modular
cyclic deformation complex $\mathrm{Def}_{\mathrm{cyc}}^{\mathrm{mod}}(\mathcal{A})$
and universal Maurer--Cartan element
$\Theta_{\mathcal{A}} \in \mathrm{MC}(\mathrm{Def}_{\mathrm{cyc}}^{\mathrm{mod}}(\mathcal{A}))$.
The complex carries an arity filtration $F^{\bullet}$ with
$F^r = \bigoplus_{s \geq r} \mathrm{Def}_{\mathrm{cyc}}^{(s)}$,
and the \emph{shadow obstruction tower} is the sequence of
finite-order truncations
\[
 \Theta_{\mathcal{A}}^{\leq 2},\;
 \Theta_{\mathcal{A}}^{\leq 3},\;
 \Theta_{\mathcal{A}}^{\leq 4},\;
 \ldots
\]
where $\Theta_{\mathcal{A}}^{\leq r}$ solves the Maurer--Cartan equation
in the quotient $\mathfrak{g}_{\mathcal{A}}^{\mathrm{amb}} / F^{r+1}$.
The arity-$r$ projection $\mathrm{Sh}_r(\mathcal{A})$ is the
\emph{$r$-th shadow} of $\mathcal{A}$. On a one-dimensional
slice of the deformation complex, $\mathrm{Sh}_r$ is determined
by a single scalar $S_r$; the MC equation projected to arity $r$
gives the \emph{all-arity master equation}
\begin{equation}
\label{eq:master-equation}
 \nabla_H(\mathrm{Sh}_r) + \mathfrak{o}^{(r)} = 0,
\end{equation}
where $\nabla_H$ is the linearized operator and $\mathfrak{o}^{(r)}$
is the obstruction from lower arities.

\subsection{Primary lines and the shadow metric}

\begin{definition}[Primary line]
\label{def:primary-line}
A \emph{primary line} in $\mathrm{Def}_{\mathrm{cyc}}^{\mathrm{mod}}(\mathcal{A})$
is a one-dimensional subspace $L = \mathbf{k} \cdot \eta$
spanned by a single primary field $\eta$ of definite conformal
weight $h_\eta$. For a chiral algebra with generators
$\phi_1, \ldots, \phi_N$ of weights $h_1, \ldots, h_N$,
each generator spans a primary line; the full first-order
deformation space is $\bigoplus_i L_i$.
\end{definition}

Let $L$ be a primary line with coordinate $x$, and write
$\mathrm{Sh}_r|_L = S_r \, x^r$ for the shadow data restricted
to $L$. Set
\[
 \kappa := S_2, \qquad \alpha := S_3.
\]
The scalar $\kappa$ is the modular characteristic (curvature)
of $\mathcal{A}$ on~$L$; the scalar $\alpha$ is the cubic shadow.

\begin{definition}[Shadow metric]
\label{def:shadow-metric-paper}
Define the \emph{weighted initial data}
$a_0 := 2\kappa$, $a_1 := 3\alpha$, $a_2 := 4S_4$.
The \emph{shadow metric} (shadow quadratic form) on $L$ is the
polynomial
\begin{equation}
\label{eq:shadow-metric}
 Q_L(t)
 \;:=\;
 a_0^2 + 2\,a_0\,a_1\,t + (a_1^2 + 2\,a_0\,a_2)\,t^2
 \;=\;
 4\kappa^2 + 12\kappa\alpha\,t + (9\alpha^2 + 16\kappa S_4)\,t^2
 \;\in\; k(c)[t].
\end{equation}
The \emph{critical discriminant} of $L$ is
\begin{equation}
\label{eq:critical-discriminant}
 \Delta := 8\kappa S_4 = a_0 \, a_2.
\end{equation}
\end{definition}

\begin{remark}
The term ``metric'' is used by analogy with the role of $Q_L$
in governing the geometry of the shadow obstruction tower, not in the sense
of a Riemannian metric. The shadow metric is a quadratic
polynomial in the arity parameter $t$, determined by three OPE
invariants: $\kappa$, $\alpha$, and $S_4$.
\end{remark}

\subsection{The recursion on a primary line}

On a primary line $L$, the all-arity master equation
\eqref{eq:master-equation} reduces to a recursion among the scalars
$S_r$.

\begin{proposition}[Single-line recursion]
\label{prop:single-line-recursion}
On a primary line $L$ with propagator $P := 1/\kappa$,
the MC equation at arity $r \geq 4$ is
\begin{equation}
\label{eq:single-line-recursion}
 S_r
 \;=\;
 -\frac{P}{2r}
 \sum_{\substack{j + k = r + 2 \\ 3 \leq j \leq k}}
 c_{jk}\, j\, k\, S_j \, S_k,
 \qquad
 c_{jk} =
 \begin{cases}
 1 & \text{if } j < k, \\
 1/2 & \text{if } j = k.
 \end{cases}
\end{equation}
This recursion is well-defined: since $j \geq 3$ and $j + k = r + 2$,
every term on the right has $k = r + 2 - j \leq r - 1$, so $S_r$
is determined by $\{S_2, S_3, \ldots, S_{r-1}\}$.
\end{proposition}

\begin{proof}
The Lie bracket on the cyclic deformation complex, restricted to
$L$, pairs $\mathrm{Sh}_j = S_j x^j$ and
$\mathrm{Sh}_k = S_k x^k$ to produce a term of arity $j + k - 2$
with coefficient $jk\,P\,S_j S_k$. The eigenvalue identity
$\nabla_H(x^r) = 2r \, x^r$ separates the MC equation at
arity $r$ into the linear term $2r \, S_r$ on the left and the
obstruction sum on the right, giving the stated recursion.
The symmetry factor $c_{jk}$ accounts for the factor
$\tfrac{1}{2}$ in $\tfrac{1}{2}[\Theta, \Theta]$ when $j = k$.
\end{proof}

\subsection{The weighted shadow generating function}

\begin{definition}[Weighted shadow generating function]
\label{def:shadow-gf}
The \emph{weighted shadow generating function} on $L$ is
\begin{equation}
\label{eq:shadow-gf}
 H(t) \;:=\; \sum_{r \geq 2} r \, S_r \, t^r.
\end{equation}
Setting $a_n := (n+2)\,S_{n+2}$ for $n \geq 0$, this becomes
$H(t) = \sum_{n \geq 0} a_n \, t^{n+2}$, with initial values
$a_0 = 2\kappa$, $a_1 = 3\alpha$, $a_2 = 4S_4$.
\end{definition}

The weighted generating function $H$ packages the recursion
\eqref{eq:single-line-recursion} as a convolution identity.
Define $F(t) := H(t)/t^2 = \sum_{n \geq 0} a_n \, t^n$.
The Cauchy product $F^2 = \sum_{m \geq 0}
(\sum_{i+j=m} a_i a_j)\, t^m$ has its coefficients at
$m \geq 3$ controlled by the MC recursion. The next theorem
shows that the MC equation forces $F^2$ to be a polynomial
of degree $2$.


% ================================================================
% §4. Algebraicity and the Depth Classification
% ================================================================

\section{Algebraicity and the depth classification}
\label{sec:algebraicity}

\subsection{The Riccati algebraicity theorem}

\begin{theorem}[Riccati algebraicity]
\label{thm:riccati-algebraicity}
Let $L$ be a primary line with shadow data $S_r$, shadow metric
$Q_L$ \textup{(}Definition~\textup{\ref{def:shadow-metric-paper}}\textup{)},
and weighted shadow generating function $H$
\textup{(}Definition~\textup{\ref{def:shadow-gf}}\textup{)}.
Then the all-arity master equation on $L$ is equivalent to the
algebraic relation
\begin{equation}
\label{eq:riccati-algebraic-relation}
 H(t)^2 \;=\; t^4 \, Q_L(t).
\end{equation}
The shadow generating function is therefore algebraic of
degree $2$ over $k(\kappa, \alpha, S_4)(t)$:
\begin{equation}
\label{eq:shadow-closed-form}
 H(t) \;=\; t^2 \sqrt{Q_L(t)}.
\end{equation}
\end{theorem}

\begin{proof}
Set $F(t) := H(t)/t^2 = \sum_{n \geq 0} a_n \, t^n$ where
$a_n = (n+2)\, S_{n+2}$.
The claim \eqref{eq:riccati-algebraic-relation} is equivalent to
$F(t)^2 = Q_L(t)$. We verify this by computing the Cauchy
product $F^2 = \sum_{m \geq 0} \bigl(\sum_{i+j=m} a_i a_j\bigr)
\, t^m$ and showing it agrees with $Q_L$ at all orders.

\emph{Orders $0$, $1$, $2$ (initial data).}
At $m = 0$: $a_0^2 = (2\kappa)^2 = 4\kappa^2$.
At $m = 1$: $2\,a_0\,a_1 = 2 \cdot 2\kappa \cdot 3\alpha = 12\kappa\alpha$.
At $m = 2$: $a_1^2 + 2\,a_0\,a_2 = 9\alpha^2 + 2 \cdot 2\kappa \cdot 4S_4
= 9\alpha^2 + 16\kappa S_4$.
These are exactly the coefficients of $Q_L(t)$ in
\eqref{eq:shadow-metric}.

\emph{Orders $m \geq 3$ (the MC recursion forces vanishing).}
Since $Q_L$ has degree $2$ in $t$, the identity $F^2 = Q_L$ requires
\begin{equation}
\label{eq:convolution-vanishing}
 \sum_{i + j = m} a_i \, a_j \;=\; 0
 \qquad \text{for all } m \geq 3.
\end{equation}
We show that this is exactly the MC recursion
\eqref{eq:single-line-recursion} at arity $r = m + 2$.

The sum $\sum_{i+j=m} a_i a_j$ with
$a_i = (i+2)\,S_{i+2}$ and $a_j = (j+2)\,S_{j+2}$
equals $\sum_{p+q = m+4,\, p,q \geq 2} p\,q\,S_p\,S_q$,
where $p = i+2$, $q = j+2$. Setting $r = m+2$, this becomes
$\sum_{p+q = r+2,\, p,q \geq 2} p\,q\,S_p\,S_q$.

Isolating the terms with $p = 2$ or $q = 2$
(which contribute $2\kappa$ from $S_2 = \kappa$):
\begin{align*}
 \sum_{p+q=r+2,\, p,q \geq 2} pq\,S_p S_q
 &= 2 \cdot 2r \cdot \kappa\,S_r
 + \sum_{\substack{p+q=r+2 \\ p,q \geq 3}} pq\,S_p S_q \\
 &= 4r\kappa\,S_r
 + 2\!\!\sum_{\substack{j+k=r+2 \\ 3 \leq j \leq k}}\!\!
 c_{jk}\,jk\,S_j S_k,
\end{align*}
where the factor $2$ in front of the restricted sum accounts for
$S_p S_q + S_q S_p$ when $p \neq q$, and $c_{jk}$ corrects
for the diagonal. Substituting the recursion
\eqref{eq:single-line-recursion}:
\[
 2\!\!\sum_{\substack{j+k=r+2 \\ 3 \leq j \leq k}}\!\!
 c_{jk}\,jk\,S_j S_k
 \;=\;
 2 \cdot (-2r\kappa\,S_r)
 \;=\;
 -4r\kappa\,S_r.
\]
Hence $\sum_{p+q=r+2} pq\,S_p S_q = 4r\kappa S_r - 4r\kappa S_r = 0$,
which is \eqref{eq:convolution-vanishing}.

Conversely, the vanishing \eqref{eq:convolution-vanishing} at
order $m \geq 3$ implies the recursion
\eqref{eq:single-line-recursion} at $r = m+2$ by solving for
$S_r$ (using $\kappa \neq 0$).

The equivalence of \eqref{eq:riccati-algebraic-relation} with
the full MC recursion is therefore established at all arities.
Taking positive square roots (the branch with $F(0) = a_0 = 2\kappa > 0$)
gives $H(t) = t^2 F(t) = t^2 \sqrt{Q_L(t)}$.
\end{proof}

\begin{remark}[Riccati ODE reformulation]
\label{rem:riccati-ode}
Set $U(t) := \sum_{r \geq 3} S_r \, t^r$ (the nonlinear part,
removing $\mathrm{Sh}_2$). The algebraic relation
\eqref{eq:riccati-algebraic-relation} is equivalent to the
Riccati-type first-order ODE
\begin{equation}
\label{eq:riccati-ode}
 2t\,U'(t) + \frac{P}{2}\bigl(U'(t)\bigr)^2
 = R_3\,t^3 + R_4\,t^4,
\end{equation}
where $R_3 = 6\alpha$ and
$R_4 = (9\alpha^2 + 2\Delta)/(2\kappa)$ encode the cubic
and quartic inputs. The left-hand side is quadratic in $U'$:
this is the precise sense in which the shadow obstruction tower is algebraic
of degree $2$.
\end{remark}

\subsection{Gaussian decomposition}

\begin{corollary}[Gaussian decomposition]
\label{cor:gaussian-decomposition}
The shadow metric decomposes as
\begin{equation}
\label{eq:gaussian-decomposition}
 Q_L(t) \;=\; (2\kappa + 3\alpha\,t)^2 + 2\Delta\,t^2,
\end{equation}
where $\Delta = 8\kappa S_4$ is the critical discriminant.
The first term $(2\kappa + 3\alpha t)^2$ is the \emph{Gaussian
envelope}: the shadow metric of the truncated tower with
$S_r = 0$ for $r \geq 4$. The second term $2\Delta\,t^2$
is the \emph{interaction correction}, measuring the quartic
OPE contribution not captured by the cubic data alone.
\end{corollary}

\begin{proof}
$(2\kappa + 3\alpha t)^2 = 4\kappa^2 + 12\kappa\alpha\,t
+ 9\alpha^2\,t^2$. Adding $2\Delta\,t^2 = 16\kappa S_4\,t^2$
recovers the coefficient $9\alpha^2 + 16\kappa S_4$ of $t^2$
in $Q_L(t)$.
\end{proof}

\subsection{Intrinsic quartic principle}

\begin{corollary}[Intrinsic quartic principle]
\label{cor:intrinsic-quartic}
On a primary line $L$, the shadow obstruction tower has intrinsic OPE data
at arities $2$, $3$, $4$ only. All shadow coefficients $S_r$
for $r \geq 5$ are determined by $(\kappa, \alpha, S_4)$
alone, via the closed form
\begin{equation}
\label{eq:shadow-from-closed-form}
 S_r \;=\; \frac{1}{r}\,[t^{r-2}]\,\sqrt{Q_L(t)}
 \qquad (r \geq 2).
\end{equation}
\end{corollary}

\begin{proof}
$Q_L(t) = 4\kappa^2 + 12\kappa\alpha\,t + (9\alpha^2 + 16\kappa S_4)\,t^2$
depends on exactly $(\kappa, \alpha, S_4)$.
Since $H(t) = t^2\sqrt{Q_L(t)} = \sum_{r \geq 2} rS_r\,t^r$
by Theorem~\ref{thm:riccati-algebraicity}, comparing
coefficients gives $rS_r = [t^{r-2}]\sqrt{Q_L}$.
\end{proof}

\subsection{The shadow connection}

\begin{proposition}[Shadow connection]
\label{prop:shadow-connection}
The shadow metric defines a logarithmic connection on
$\mathcal{O}_L$ over $L \setminus \{Q_L = 0\}$:
\begin{equation}
\label{eq:shadow-connection}
 \nabla^{\mathrm{sh}} \;=\; d - \omega,
 \qquad
 \omega \;:=\; \frac{Q_L'(t)}{2\,Q_L(t)}\,dt
 \;=\; d\bigl(\log\sqrt{Q_L}\bigr).
\end{equation}
\begin{enumerate}[label=\textup{(\roman*)}]
\item The connection is flat: $(\nabla^{\mathrm{sh}})^2 = 0$.
\item At each simple zero $t_0$ of $Q_L$, the residue
 of $\omega$ is $\tfrac{1}{2}$.
\item The monodromy is $\exp(2\pi i \cdot \tfrac{1}{2}) = -1$:
 the standard monodromy of a square root around a simple
 branch point.
\item The flat section with $\Phi(0) = 1$ is
 $\Phi(t) = \sqrt{Q_L(t)/Q_L(0)} = \sqrt{Q_L(t)}/(2\kappa)$,
 and the shadow generating function is
 $H(t) = 2\kappa\,t^2\,\Phi(t)$: the parallel transport of
 the curvature along the arity parameter, weighted by $t^2$.
\end{enumerate}
\end{proposition}

\begin{proof}
Flatness is automatic: $\omega$ is a closed $1$-form on a
$1$-dimensional base. At a simple zero $t_0$ of $Q_L$:
$Q_L(t) \approx Q_L'(t_0)(t - t_0)$, so
$\omega \approx \frac{1}{2(t - t_0)}\,dt$,
giving residue $\tfrac{1}{2}$. The monodromy exponentiation
is immediate. For (iv): $\nabla^{\mathrm{sh}}\Phi = 0$
since $d\Phi = \omega\,\Phi$ by construction, and the
evaluation at $t = 0$ gives
$\Phi(0) = \sqrt{Q_L(0)}/(2\kappa) = 2\kappa/(2\kappa) = 1$.
\end{proof}

\begin{remark}
The shadow generating function $H(t) = 2\kappa\,t^2\,\Phi(t)$
is \emph{not} a flat section of $\nabla^{\mathrm{sh}}$: the
factor $t^2$ contributes $\nabla^{\mathrm{sh}}(H) \neq 0$.
The flat section is $\Phi(t) = \sqrt{Q_L(t)}/(2\kappa)$; the
shadow obstruction tower is the Taylor expansion of $t^2 \cdot \Phi(t)$,
reflecting the arity offset (the tower starts at arity $2$,
not arity $0$).
\end{remark}

\subsection{The single-line dichotomy}

\begin{theorem}[Single-line dichotomy]
\label{thm:single-line-dichotomy}
Let $L$, $S_r$, $\alpha$, $\Delta$, $Q_L$ be as above,
and define the shadow depth $r_{\max}|_L := \max\{r : S_r \neq 0\}$
\textup{(}with $r_{\max} = \infty$ if $S_r \neq 0$ for
infinitely many $r$\textup{)}.
Then $r_{\max}|_L \in \{2, 3, \infty\}$, classified by the
shadow metric:
\begin{enumerate}[label=\textup{(\roman*)}]
\item \textbf{Class $\mathbf{G}$}
 \textup{(}$\Delta = 0$, $\alpha = 0$\textup{)}.\enspace
 $Q_L = (2\kappa)^2$ is constant.
 $H(t) = 2\kappa\,t^2$. Only $S_2$ survives: $r_{\max} = 2$.

\item \textbf{Class $\mathbf{L}$}
 \textup{(}$\Delta = 0$, $\alpha \neq 0$\textup{)}.\enspace
 $Q_L = (2\kappa + 3\alpha\,t)^2$ is a perfect square.
 $H(t) = t^2(2\kappa + 3\alpha\,t)$. $S_2$ and $S_3$
 survive: $r_{\max} = 3$.

\item \textbf{Class $\mathbf{M}$}
 \textup{(}$\Delta \neq 0$\textup{)}.\enspace
 $Q_L$ is irreducible over $k(c)$; $\sqrt{Q_L}$ is
 irrational; $S_r \neq 0$ for all $r \geq 4$ generically:
 $r_{\max} = \infty$.
\end{enumerate}
For $r \geq 4$,
\begin{equation}
\label{eq:universal-factorization}
 S_r = \Delta \cdot R_r,
 \qquad R_r \in \mathbb{Q}(\alpha, \kappa, S_4).
\end{equation}
\end{theorem}

\begin{proof}
The classification follows from Theorem~\ref{thm:riccati-algebraicity}
and the Gaussian decomposition
(Corollary~\ref{cor:gaussian-decomposition}).

\emph{Case $\Delta = 0$.}
Then $Q_L = (2\kappa + 3\alpha t)^2$ is a perfect square, so
$\sqrt{Q_L} = 2\kappa + 3\alpha t$ is polynomial and
$H(t) = t^2(2\kappa + 3\alpha t)$ terminates.
If $\alpha = 0$: $H(t) = 2\kappa t^2$, giving $r_{\max} = 2$.
If $\alpha \neq 0$: $H(t) = 2\kappa t^2 + 3\alpha t^3$,
giving $r_{\max} = 3$.

\emph{Case $\Delta \neq 0$.}
The discriminant of $Q_L$ as a polynomial in $t$ is
\[
 (12\kappa\alpha)^2 - 4 \cdot 4\kappa^2 \cdot
 (9\alpha^2 + 16\kappa S_4) = -32\kappa^2\Delta \neq 0,
\]
so $Q_L$ is irreducible over $k(c)$. The Taylor expansion
of $\sqrt{Q_L}$ around $t = 0$ has nonzero coefficients at
all orders generically, hence $r_{\max} = \infty$.

\emph{Universal factorization.}
The binomial expansion of $\sqrt{Q_L}$ around the Gaussian
envelope $(2\kappa + 3\alpha t)$ is
\[
 \sqrt{Q_L} = (2\kappa + 3\alpha t)\sqrt{1 + \frac{2\Delta\,t^2}
 {(2\kappa + 3\alpha t)^2}}.
\]
Expanding $\sqrt{1 + u}$ as $1 + \tfrac{1}{2}u - \tfrac{1}{8}u^2
+ \cdots$ with $u = 2\Delta t^2/(2\kappa + 3\alpha t)^2$,
every coefficient of $t^n$ for $n \geq 2$ in
$\sqrt{Q_L} - (2\kappa + 3\alpha t)$ carries a factor of
$\Delta$. Since $rS_r = [t^{r-2}]\sqrt{Q_L}$ for $r \geq 2$,
and the Gaussian part $2\kappa + 3\alpha t$ contributes only
at $r \leq 3$, the factorization $S_r = \Delta \cdot R_r$
for $r \geq 4$ follows.
\end{proof}

\begin{remark}[The contact class]
\label{rem:contact-class}
The partition $\mathbf{G}/\mathbf{L}/\mathbf{C}/\mathbf{M}$
of modular Koszul algebras includes a fourth class $\mathbf{C}$
(contact, $r_{\max} = 4$), realized by the $\beta\gamma$ system.
This class escapes the single-line dichotomy by \emph{stratum
separation}: the quartic contact invariant lives on a charged
stratum of $\mathrm{Def}_{\mathrm{cyc}}^{\mathrm{mod}}$ where
$\kappa = 0$, and the shadow metric requires $\kappa \neq 0$.
On the charged stratum, the self-bracket maps to a zero-dimensional
target by rank-one rigidity, so the quartic pump does not activate
and the tower terminates at $r = 4$. The algebraic shadow depth
$d_{\mathrm{alg}} \in \{0, 1, 2, \infty\}$ exhausts the four
classes: $\mathbf{G}$ ($d_{\mathrm{alg}} = 0$),
$\mathbf{L}$ ($d_{\mathrm{alg}} = 1$),
$\mathbf{C}$ ($d_{\mathrm{alg}} = 2$),
$\mathbf{M}$ ($d_{\mathrm{alg}} = \infty$).
\end{remark}

\subsection{The shadow growth rate}

For class $\mathbf{M}$ algebras, the shadow obstruction tower is infinite
but governed by a single analytic invariant: the growth rate of
the coefficients $S_r$.

\begin{definition}[Shadow growth rate]
\label{def:shadow-growth-rate}
Let $L$ be a primary line of class $\mathbf{M}$ ($\Delta \neq 0$),
and write $Q_L(t) = q_0 + q_1 t + q_2 t^2$ for the shadow metric.
The \emph{shadow growth rate} of $L$ is
\begin{equation}
\label{eq:shadow-growth-rate}
 \rho_L \;:=\; \sqrt{\frac{q_2}{q_0}}
 \;=\;
 \frac{\sqrt{9\alpha^2 + 2\Delta}}{2|\kappa|}.
\end{equation}
By Vieta's theorem, the zeros $t_0$, $\bar{t}_0$ of $Q_L$
satisfy $|t_0|^2 = q_0/q_2$, so
$\rho_L = 1/|t_0|$: the reciprocal of the branch-point modulus.
For classes $\mathbf{G}$ and $\mathbf{L}$, set $\rho := 0$
by convention.
\end{definition}

\begin{theorem}[Shadow asymptotics]
\label{thm:shadow-asymptotics}
Let $L$ be a primary line of class $\mathbf{M}$ with shadow
growth rate $\rho_L$.
\begin{enumerate}[label=\textup{(\roman*)}]
\item \emph{Asymptotic formula.}
 \begin{equation}
 \label{eq:shadow-asymptotic-formula}
 S_r \;\sim\; A(\mathcal{A})\,\rho_L^{\,r}\,r^{-5/2}\,
 \cos(r\theta + \varphi)
 \qquad (r \to \infty),
 \end{equation}
 where $\theta := -\arg(t_0)$ is the oscillation phase,
 $A(\mathcal{A}) > 0$, and $\varphi$ is computable from the
 branch-point residue of $\sqrt{Q_L}$. The exponent $-5/2$
 arises from $S_r = [t^{r-2}]\sqrt{Q_L}/r$ combined with the
 $n^{-3/2}$ transfer law for square-root singularities.

\item \emph{Unit circle criterion.}
 $\rho_L = 1$ if and only if $q_0 = q_2$, i.e.,
 \begin{equation}
 \label{eq:unit-circle-criterion}
 (2\kappa - 3\alpha)(2\kappa + 3\alpha) = 2\Delta.
 \end{equation}
 Equivalently, the branch points of $Q_L$ lie on the unit
 circle $|t| = 1$.

\item \emph{Convergence partition.}
 $\rho_L < 1$: the shadow obstruction tower converges
 absolutely.
 $\rho_L > 1$: the shadow obstruction tower diverges, but the algebraic
 function $\sqrt{Q_L(t)}$ extends analytically past the
 convergence radius.
 $\rho_L = 1$: marginal case (the critical surface).
\end{enumerate}
\end{theorem}

\begin{proof}
The function $\sqrt{Q_L(t)}$ has algebraic branch-cut
singularities of order $\tfrac{1}{2}$ at each simple zero
$t_0$ of $Q_L$. Since $Q_L$ is quadratic with nonzero
discriminant ($\Delta \neq 0$ implies
$\mathrm{disc}(Q_L) = -32\kappa^2\Delta \neq 0$),
the two zeros $t_0, \bar{t}_0$ are distinct with
$|t_0| = |{\bar t}_0| = \sqrt{q_0/q_2} = 1/\rho_L$.

The transfer theorem for algebraic singularities
(Flajolet--Sedgewick, Theorem~VI.1) gives
$[t^n]\sqrt{Q_L} \sim C_0\,t_0^{-n}\,n^{-3/2}$
for each dominant singularity $t_0$ at distance
$R = 1/\rho_L$ from the origin. Both branch points
contribute (equal modulus), producing the cosine
interference term from
$2\operatorname{Re}(C_0\,t_0^{-(r-2)})$.
Since $S_r = [t^{r-2}]\sqrt{Q_L}/r$
(Corollary~\ref{cor:intrinsic-quartic}), the extra
factor $1/r$ gives the stated $r^{-5/2}$ exponent.

For (ii): $\rho_L = 1$ iff $q_2/q_0 = 1$. Writing
$q_0 = 4\kappa^2$ and $q_2 = 9\alpha^2 + 2\Delta$, the
condition $q_0 = q_2$ factors as
$4\kappa^2 - 9\alpha^2 = 2\Delta$.
Part (iii) is immediate from the definition of the
convergence radius $R = 1/\rho_L$.
\end{proof}

\subsection{The depth decomposition}

\begin{theorem}[Depth decomposition]
\label{thm:depth-decomposition}
The shadow depth of a modular Koszul algebra $\mathcal{A}$
decomposes as
\begin{equation}
\label{eq:depth-decomposition}
 d(\mathcal{A}) \;=\; 1 + d_{\mathrm{arith}}(\mathcal{A})
 + d_{\mathrm{alg}}(\mathcal{A}),
\end{equation}
where:
\begin{enumerate}[label=\textup{(\roman*)}]
\item $d_{\mathrm{arith}}$ counts the independent holomorphic
 Hecke eigenforms in the spectral decomposition of the completed
 shadow partition function $\widehat{Z}^c_{\mathcal{A}}$
 on $\mathcal{M}_{1,1}$;

\item $d_{\mathrm{alg}} \in \{0, 1, 2, \infty\}$ is the
 \emph{algebraic depth}, determined by the single-line dichotomy
 \textup{(}Theorem~\textup{\ref{thm:single-line-dichotomy}}\textup{)}
 and stratum separation
 \textup{(}Remark~\textup{\ref{rem:contact-class}}\textup{)}:
 \[
 d_{\mathrm{alg}} =
 \begin{cases}
 0 & \text{class } \mathbf{G}, \\
 1 & \text{class } \mathbf{L}, \\
 2 & \text{class } \mathbf{C}, \\
 \infty & \text{class } \mathbf{M};
 \end{cases}
 \]

\item depths $d \geq 5$ are purely arithmetic: they arise from
 $d_{\mathrm{arith}}$, not from the $A_\infty$ structure.
\end{enumerate}
\end{theorem}

\begin{proof}
The algebraic depth $d_{\mathrm{alg}}$ is determined by the
structure of the cyclic deformation complex and the MC recursion:
the single-line dichotomy
(Theorem~\ref{thm:single-line-dichotomy}) gives
$d_{\mathrm{alg}} \in \{0, 1, \infty\}$ on each primary line,
and stratum separation (Remark~\ref{rem:contact-class}) provides
$d_{\mathrm{alg}} = 2$ for the contact class.
Since $d_{\mathrm{alg}} \leq 2$ for finite cases
($\mathbf{G}/\mathbf{L}/\mathbf{C}$) and the base term
is $1$ (the curvature $\kappa$ at arity $2$), the maximum
finite depth without arithmetic contributions is
$1 + 0 + 2 = 3$; hence all depths $d \geq 5$ require
$d_{\mathrm{arith}} \geq 2$.

The arithmetic depth counts independent holomorphic Hecke
eigenforms by decomposing $\widehat{Z}^c_{\mathcal{A}}$
via the Roelcke--Selberg spectral decomposition on
$\mathcal{M}_{1,1}$. Each holomorphic eigenform of weight $k$
contributes a critical line at $\operatorname{Re}(s) = k/2$;
Maass cusp-form projections contribute no new critical lines.
\end{proof}

\begin{example}[Standard families]
\label{ex:standard-families}
The shadow metric data for the standard families:
\begin{center}
\small
\renewcommand{\arraystretch}{1.3}
\begin{tabular}{lcccccl}
\toprule
Family & $\kappa$ & $\alpha$ & $S_4$ & $\Delta$ & Class
 & $\rho$ \\
\midrule
Heisenberg $\mathcal{H}_k$ & $k$ & $0$ & $0$ & $0$
 & $\mathbf{G}$ & $0$ \\
Affine $\widehat{\mathfrak{g}}_k$
 & $\frac{d(k+h^\vee)}{2h^\vee}$ & $\neq 0$ & $0$
 & $0$ & $\mathbf{L}$ & $0$ \\
$\beta\gamma$ (weight $\lambda$) & $6\lambda^2{-}6\lambda{+}1$ & $0$ & $\neq 0$ & $0^*$
 & $\mathbf{C}$ & $*$ \\
Virasoro $\mathrm{Vir}_c$ & $c/2$ & $2$
 & $\frac{10}{c(5c+22)}$ & $\frac{40}{5c+22}$
 & $\mathbf{M}$ & $\frac{1}{c}\sqrt{\frac{180c+872}{5c+22}}$ \\
\bottomrule
\end{tabular}
\end{center}
\vspace{2pt}
\noindent
${}^*$\,Stratum separation: $\Delta = 0$ on the primary line
(where $S_4 = 0$); the quartic contact invariant
$Q^{\mathrm{contact}} \neq 0$ lives on a charged stratum
where $\kappa = 0$, and the single-line dichotomy does not
apply.

For the Virasoro algebra, the shadow metric evaluates to
\[
 Q_{\mathrm{Vir}}(t) = c^2 + 12c\,t
 + \frac{180c + 872}{5c + 22}\,t^2,
\]
with Gaussian decomposition
$Q_{\mathrm{Vir}} = (c + 6t)^2 + \frac{80}{5c+22}\,t^2$.
The critical central charge $c^*$ where $\rho = 1$ is the
unique positive real root of
\[
 5c^3 + 22c^2 - 180c - 872 = 0,
\]
approximately $c^* \approx 6.12$. For $c > c^*$ the shadow
tower converges; for $c < c^*$ it diverges. At the Koszul
self-dual point $c = 13$: $\rho(\mathrm{Vir}_{13}) \approx 0.467$.
\end{example}

\begin{remark}[Spectral curve]
The algebraic relation $H^2 = t^4 Q_L(t)$ defines, for each
primary line $L$, a spectral curve
$\Sigma_L := \{H^2 = t^4 Q_L(t)\} \subset \mathbb{A}^2_{t,H}$.
Since $Q_L$ is quadratic in $t$, $\Sigma_L$ has arithmetic
genus $0$; the projection $(t,H) \mapsto t$ is a double cover
ramified at the zeros of $Q_L$. The degeneration loci in
parameter space are controlled by the critical discriminant:
$\Sigma_L$ acquires a node precisely when $\Delta = 0$.
The shadow growth rate $\rho_L$ is the reciprocal of the
distance from $t = 0$ to the nearest ramification point.
\end{remark}


%% classification.tex — §5–6 of the standalone paper
%% The G/L/C/M classification and shadow tables

%% =====================================================================
%%  §5.  The Four Shadow Classes
%% =====================================================================

\section{The Four Shadow Classes}\label{sec:four-classes}

The single-line dichotomy (Theorem~\ref{thm:single-line-dichotomy})
partitions the space of modular Koszul algebras into four shadow classes,
determined by two invariants of the shadow metric
$Q_L(t) = (2\kappa + 3\alpha t)^2 + 2\Delta\, t^2$:
the cubic coefficient~$\alpha = S_3$ and the critical
discriminant $\Delta = 8\kappa S_4$.

\begin{definition}[Shadow classes]\label{def:shadow-classes}
Let $\cA$ be a modular Koszul algebra with shadow metric $Q_L$ on a
one-dimensional primary line $L$.  The \emph{shadow class} of $\cA$
along $L$ is:
\[
\renewcommand{\arraystretch}{1.4}
\begin{array}{c|cc|cc}
\textbf{Class} & \alpha & \Delta & r_{\max} & Q_L \\
\hline
G\;\text{(Gaussian)} & 0 & 0 & 2 & (2\kappa)^2 \\
L\;\text{(Lie/tree)} & \neq 0 & 0 & 3 & (2\kappa + 3\alpha t)^2 \\
C\;\text{(Contact)} & 0 & 0^* & 4 & (2\kappa)^2\;\text{(primary line)} \\
M\;\text{(Mixed)} & \neq 0 & \neq 0 & \infty & \text{irreducible quadratic}
\end{array}
\]
Here $r_{\max}$ is the \emph{shadow depth}: the largest arity~$r$
for which the shadow coefficient $S_r \neq 0$.
${}^*$Class~C has $\Delta = 0$ on the primary line (where
the single-line dichotomy applies), but the quartic contact
invariant $Q^{\mathrm{contact}} \neq 0$ lives on a charged
stratum; see Definition~\ref{def:class-C}.
\end{definition}

The classification is exhaustive: every entry in Table~\ref{tab:master-shadow}
falls into exactly one class, and every class is realized by at least
one standard family.  The mechanism is algebraic: the shadow generating
function $H(t) = t^2\sqrt{Q_L(t)}$ is algebraic of degree~$2$
(Theorem~\ref{thm:riccati-algebraicity}), and its Taylor coefficients
$S_r$ are computed from the convolution recursion
$H(t)^2 = t^4 Q_L(t)$, where $H(t) = \sum_{r \geq 2} r\, S_r\, t^r$
and $S_r = a_{r-2}/r$ with $a_n$ the Taylor coefficients of $\sqrt{Q_L}$.
When $Q_L$ is a perfect square ($\Delta = 0$), the square root
terminates; when $Q_L$ is irreducible ($\Delta \neq 0$), the Taylor
expansion is infinite.

\medskip

\subsection{Class G: Gaussian ($r_{\max} = 2$)}

\begin{definition}\label{def:class-G}
A modular Koszul algebra $\cA$ is \emph{Gaussian} if
$S_r(\cA) = 0$ for all $r \geq 3$.  Equivalently,
$\alpha = 0$ and $\Delta = 0$, so the shadow metric
$Q_L(t) = 4\kappa^2$ is constant and the shadow tower
terminates at the scalar curvature.
\end{definition}

\noindent\textbf{Examples.}  Heisenberg $\cH_k$ with $\kappa = k$;
lattice VOAs $V_\Lambda$ with $\kappa = \operatorname{rank}(\Lambda)$;
free fermion $\psi$ with $\kappa = 1/4$;
$bc$ ghosts with $\kappa = -(6\lambda^2 - 6\lambda + 1)$.

\noindent\textbf{Characterization.}
Class~G algebras are freely generated as chiral algebras.  Their bar
complex is quadratic (the differential has no higher-order terms),
and the OPE on each primary line is abelian.  The full obstruction
is captured by a single scalar:
\[
\Theta_\cA = \kappa(\cA) \cdot \eta \otimes \Lambda,
\]
where $\eta$ is the cyclic direction and $\Lambda \in H^*(\bar{\cM}_g)$
is the tautological class.  All genus-$g$ data is controlled by
$F_g(\cA) = \kappa(\cA) \cdot \lambda_g^{\mathrm{FP}}$.

\medskip

\subsection{Class L: Lie/tree ($r_{\max} = 3$)}

\begin{definition}\label{def:class-L}
A modular Koszul algebra $\cA$ is \emph{Lie/tree} if
$\alpha \neq 0$ and $\Delta = 0$, so that $S_3 \neq 0$
but $S_r = 0$ for all $r \geq 4$.  The shadow metric is
a perfect square:
$Q_L(t) = (2\kappa + 3\alpha t)^2$.
\end{definition}

\noindent\textbf{Examples.}  Affine Kac--Moody $V_k(\fg)$ for every
simple Lie algebra $\fg$ and every non-critical level
$k \neq -h^\vee$: type~$A$ ($\mathfrak{sl}_N$),
types $B$, $C$, $D$ (orthogonal and symplectic),
exceptional types $E_6$, $E_7$, $E_8$, $G_2$, $F_4$.
For all of these:
\[
\kappa\bigl(V_k(\fg)\bigr)
= \frac{\dim(\fg)\,(k + h^\vee)}{2h^\vee}\,.
\]

\noindent\textbf{Characterization.}
The nonzero cubic shadow $S_3 = \alpha$ arises from the Lie bracket
(structure constants $f_{abc}$), while the Jacobi identity forces the
quartic shadow $S_4 = 0$ on the primary line.  The $r$-matrix
$r(z) = \Omega_\fg/z$ (single simple pole; the bar construction absorbs
one power of the OPE pole) encodes all nonlinear structure.  The tree-level
graph sum terminates at trivalent vertices.

\medskip

\subsection{Class C: Contact ($r_{\max} = 4$)}

\begin{definition}\label{def:class-C}
A modular Koszul algebra $\cA$ is \emph{contact} if
$\Delta = 0$ on every primary line (so the single-line
dichotomy gives $r_{\max}|_L \leq 3$ on each line), but there
exists a charged stratum direction with $S_4 \neq 0$
(the quartic contact invariant), and the tower terminates at
arity~$4$: $S_r = 0$ for $r \geq 5$.
\end{definition}

\noindent\textbf{Mechanism.}
Class~C escapes the single-line dichotomy (which gives $r_{\max} \in
\{2, 3, \infty\}$) via \emph{stratum separation}: the quartic contact
invariant $Q^{\mathrm{contact}}$ lives on a charged stratum whose
self-bracket exits the cyclic deformation complex at arity~$5$, by
rank-one rigidity.  On the weight-changing primary line, the tower is
abelian ($\alpha = 0$, $S_4 = 0$); the quartic contact is visible
only on the full two-dimensional deformation space.

\noindent\textbf{Examples.}
The $\beta\gamma$ system at conformal weight~$\lambda$, with
\[
\kappa(\beta\gamma_\lambda)
= 6\lambda^2 - 6\lambda + 1
= \tfrac{1}{2}\,c(\beta\gamma_\lambda),
\]
and its Koszul dual, the $bc$ system at the same weight.
At the standard weight $\lambda = 0$ or $\lambda = 1$:
$c = 2$ and $\kappa = 1$.

\medskip

\subsection{Class M: Mixed ($r_{\max} = \infty$)}

\begin{definition}\label{def:class-M}
A modular Koszul algebra $\cA$ is \emph{mixed} if
$\alpha \neq 0$ and $\Delta \neq 0$, so the shadow metric
$Q_L$ is an irreducible quadratic in~$t$ and the shadow tower
does not terminate.
\end{definition}

\noindent\textbf{Examples.}
Virasoro $\mathrm{Vir}_c$ with $\kappa = c/2$;
$\cW$-algebras $\cW_N$ with $\kappa = (H_N - 1)\,c$ where
$H_N = \sum_{j=1}^{N} 1/j$;
Drinfeld--Sokolov reductions of affine Kac--Moody algebras.

\noindent\textbf{Characterization.}
The shadow growth rate
\[
\rho(\cA) = \frac{\sqrt{9\alpha^2 + 2\Delta}}{2|\kappa|}
\]
controls the asymptotic behaviour of $S_r$: the coefficients satisfy
$S_r \sim C(\cA)\, \rho(\cA)^r\, r^{-5/2}\, \cos(r\theta + \phi)$
for a branch-point argument~$\theta$ determined by the shadow metric.
The critical central charge $c^*$ where $\rho(\mathrm{Vir}_c) = 1$ is the
unique positive real root of the cubic
\begin{equation}\label{eq:critical-cubic}
5c^3 + 22c^2 - 180c - 872 = 0,
\qquad
c^* \approx 6.125\,.
\end{equation}
For $c > c^*$ the tower converges ($\rho < 1$), and for $c < c^*$ it
diverges ($\rho > 1$).  At the self-dual point $c = 13$:
$\rho(\mathrm{Vir}_{13}) \approx 0.467$.

The Virasoro shadow tower has explicitly computed coefficients through
arity~$10$:
\begin{align}
S_2 &= \tfrac{c}{2}\,, &
S_3 &= 2\,, &
S_4 &= \frac{10}{c(5c+22)}\,,
\label{eq:vir-S234} \\
S_5 &= \frac{-48}{c^2(5c+22)}\,, &
S_6 &= \frac{80(45c+193)}{3c^3(5c+22)^2}\,, &
S_7 &= \frac{-2880(15c+61)}{7c^4(5c+22)^2}\,.
\label{eq:vir-S567}
\end{align}

\medskip

\begin{remark}[$\mathcal{N}=2$ superconformal algebra]
\label{rem:n2-sca}
The $\mathcal{N}=2$ superconformal algebra at level~$k$ has
central charge $c = 3k/(k+2)$ and modular characteristic
$\kappa = (k+4)/4$.  The Kazama--Suzuki coset realization
gives $\kappa = \kappa(\widehat{\mathfrak{sl}}_2 \text{ at }k)
+ \kappa(\text{fermions}) - \kappa(U(1))
= 3(k{+}2)/4 + 1/2 - (k/2{+}1) = (k+4)/4$.
Under Koszul duality, $c + c' = 6$ (from the $\mathfrak{sl}_2$
involution $k \mapsto -k - 4$).  The algebra is class~$\mathbf{M}$
along the stress-tensor primary line, since the Virasoro subalgebra
has $\alpha = 2$ and $\Delta \neq 0$.
\end{remark}

\begin{remark}[Depth and Koszulness]\label{rem:depth-vs-koszulness}
Shadow depth $r_{\max}$ does \emph{not} characterize chiral Koszulness.
All four classes---Gaussian ($r_{\max} = 2$), Lie ($r_{\max} = 3$),
contact ($r_{\max} = 4$), and mixed ($r_{\max} = \infty$)---consist
of chirally Koszul algebras.  Shadow depth classifies \emph{complexity}
within the Koszul world: the operadic complexity conjecture
identifies $r_{\max}(\cA)$ with the $A_\infty$-depth and
the $L_\infty$-formality level.
\end{remark}

\begin{remark}[Galois splitting]\label{rem:galois-splitting}
The four-class partition has a Galois-theoretic interpretation.
The shadow metric $Q_L(t) = 4\kappa^2 + 12\kappa\alpha\,t
+ (9\alpha^2 + 2\Delta)\,t^2$ has discriminant
$\operatorname{disc}(Q_L) = -32\kappa^2\Delta$.
Over the splitting field:
\begin{itemize}
\item Class~$\mathbf{G}$: $Q_L = (2\kappa)^2$ is a square.
  The ``splitting field'' is $\mathbb{Q}$ itself.
  No extension; trivial Galois group.
\item Class~$\mathbf{L}$: $Q_L = (2\kappa + 3\alpha t)^2$ is
  a square with nontrivial linear factor.  The splitting field
  is still~$\mathbb{Q}$.  The single zero $t_0 = -2\kappa/(3\alpha)$
  is rational.
\item Class~$\mathbf{M}$: $\operatorname{disc}(Q_L) < 0$ at
  positive~$c$.  The zeros of $Q_L$ are a conjugate pair
  $t_0, \bar{t}_0$ in $K_L = \mathbb{Q}(\sqrt{\operatorname{disc}})$,
  an imaginary quadratic extension.  The Galois group
  $\operatorname{Gal}(K_L/\mathbb{Q}) = \mathbb{Z}/2$ acts by
  $t_0 \mapsto \bar{t}_0$, and the shadow tower inherits the
  oscillatory interference from the two conjugate branch points.
\end{itemize}
The passage from class~$\mathbf{L}$ to class~$\mathbf{M}$
is a field extension: the shadow field enlarges from $\mathbb{Q}$
to an imaginary quadratic $K_L$.  The critical discriminant
$\Delta$ is the obstruction to the extension being trivial.
\end{remark}

\begin{remark}[Total depth decomposition]\label{rem:total-depth}
The total shadow depth decomposes as
$d(\cA) = 1 + d_{\mathrm{arith}} + d_{\mathrm{alg}}$
(Theorem~\ref{thm:depth-decomposition}), where the algebraic depth
$d_{\mathrm{alg}} \in \{0, 1, 2, \infty\}$ and depths $\geq 5$ are
purely arithmetic (cusp forms contribute beyond the algebraic engine).
\end{remark}


%% =====================================================================
%%  §6.  Shadow Tables
%% =====================================================================

\section{Shadow Tables}\label{sec:shadow-tables}

All entries in the tables below are verified against the compute layer.
The modular characteristic $\kappa(\cA)$ is the arity-$2$ projection
of the universal MC element $\Theta_\cA$, and the genus-$g$ free energy
is $F_g(\cA) = \kappa(\cA) \cdot \lambda_g^{\mathrm{FP}}$, where
$\lambda_g^{\mathrm{FP}} = \frac{2^{2g-1}-1}{2^{2g-1}}
\cdot \frac{|B_{2g}|}{(2g)!}$
are the Faber--Pandharipande intersection numbers.

\subsection{Master shadow table}\label{ssec:master-shadow-table}

Table~\ref{tab:master-shadow} records, for each standard family,
the modular characteristic $\kappa$, the shadow class, the shadow
depth $r_{\max}$, the cubic coefficient $\alpha = S_3$, the quartic
coefficient $S_4$, and the critical discriminant $\Delta = 8\kappa S_4$.

\begin{table}[ht]
\centering
\caption{Master shadow table: standard landscape}
\label{tab:master-shadow}
\renewcommand{\arraystretch}{1.35}
{\small
\begin{tabular}{lcccccc}
\toprule
\textbf{Algebra $\cA$}
  & $\boldsymbol{\kappa(\cA)}$
  & \textbf{Class}
  & $\boldsymbol{r_{\max}}$
  & $\boldsymbol{S_3}$
  & $\boldsymbol{S_4}$
  & $\boldsymbol{\Delta}$ \\
\midrule
\multicolumn{7}{l}{\textit{Class G (Gaussian)}} \\[2pt]
Heisenberg $\cH_k$
  & $k$
  & G & $2$ & $0$ & $0$ & $0$ \\
Lattice $V_\Lambda$
  & $\operatorname{rank}(\Lambda)$
  & G & $2$ & $0$ & $0$ & $0$ \\
Free fermion $\psi$
  & $1/4$
  & G & $2$ & $0$ & $0$ & $0$ \\
$bc$ ghosts (weight $\lambda$)
  & $-(6\lambda^2{-}6\lambda{+}1)$
  & G & $2$ & $0$ & $0$ & $0$ \\[3pt]
\midrule
\multicolumn{7}{l}{\textit{Class L (Lie/tree)}} \\[2pt]
$\widehat{\mathfrak{sl}}_2$ at level $k$
  & $3(k{+}2)/4$
  & L & $3$
  & $\neq 0$ & $0$ & $0$ \\
$\widehat{\mathfrak{sl}}_N$ at level $k$
  & $(N^2{-}1)(k{+}N)/(2N)$
  & L & $3$
  & $\neq 0$ & $0$ & $0$ \\
$\widehat{E}_6$ at level $k$
  & $13(k{+}12)/4$
  & L & $3$
  & $\neq 0$ & $0$ & $0$ \\
$\widehat{E}_7$ at level $k$
  & $133(k{+}18)/36$
  & L & $3$
  & $\neq 0$ & $0$ & $0$ \\
$\widehat{E}_8$ at level $k$
  & $62(k{+}30)/15$
  & L & $3$
  & $\neq 0$ & $0$ & $0$ \\[3pt]
\midrule
\multicolumn{7}{l}{\textit{Class C (Contact)}} \\[2pt]
$\beta\gamma$ (weight $\lambda$)
  & $6\lambda^2{-}6\lambda{+}1$
  & C & $4$
  & $0$ & $0^\dagger$ & $0^\dagger$ \\[3pt]
\midrule
\multicolumn{7}{l}{\textit{Class M (Mixed)}} \\[2pt]
$\mathrm{Vir}_c$
  & $c/2$
  & M & $\infty$
  & $2$
  & $\dfrac{10}{c(5c{+}22)}$
  & $\dfrac{40}{5c{+}22}$ \\[8pt]
$\cW_3$ at $c$
  & $5c/6$
  & M & $\infty$
  & $\neq 0$ & $\neq 0$ & $\neq 0$ \\
$\cW_4$ at $c$
  & $13c/12$
  & M & $\infty$
  & $\neq 0$ & $\neq 0$ & $\neq 0$ \\
$\cW_5$ at $c$
  & $77c/60$
  & M & $\infty$
  & $\neq 0$ & $\neq 0$ & $\neq 0$ \\
$\cW_N$ at $c$ (general)
  & $(H_N{-}1)\,c$
  & M & $\infty$
  & $\neq 0$ & $\neq 0$ & $\neq 0$ \\
\bottomrule
\end{tabular}
}
\end{table}

\noindent
${}^\dagger$On the primary line, $S_4 = 0$ and $\Delta = 0$
for $\beta\gamma$; the quartic contact invariant
$Q^{\mathrm{contact}} \neq 0$ lives on a charged stratum
(stratum separation, Definition~\ref{def:class-C}).

\noindent
In the table, $H_N = 1 + 1/2 + \dotsb + 1/N$ denotes the
$N$-th harmonic number.  The Virasoro quartic contact
$S_4 = Q^{\mathrm{contact}}_{\mathrm{Vir}} = 10/[c(5c+22)]$
is the unique quartic shadow on the $T$-line, proved by direct
computation from the Virasoro OPE.  The T-line shadow data ($S_3 = 2$,
$S_4 = 10/[c(5c+22)]$) is \emph{universal}: it coincides for $\cW_N$
at all $N \geq 2$, since every $\cW_N$ contains the Virasoro
subalgebra along the stress-tensor primary line.  The additional
primary lines (spin-$3$, spin-$4$, etc.) carry independent shadow
data with separate quartic coefficients.


\subsection{Exceptional and non-simply-laced types}\label{ssec:exceptional-table}

Table~\ref{tab:exceptional-shadow} records the affine Kac--Moody
data for exceptional and non-simply-laced types.  All are class~L.

\begin{table}[ht]
\centering
\caption{Shadow data for exceptional and non-simply-laced types}
\label{tab:exceptional-shadow}
\renewcommand{\arraystretch}{1.35}
{\small
\begin{tabular}{lcccccc}
\toprule
$\fg$
  & $\dim(\fg)$
  & $h^\vee$
  & $\boldsymbol{\kappa\bigl(V_k(\fg)\bigr)}$
  & $c(V_k) + c(V_{k'})$
  & \textbf{Class}
  & $r_{\max}$ \\
\midrule
$B_2 \cong \mathfrak{so}_5$
  & $10$ & $3$ & $5(k{+}3)/3$ & $20$ & L & $3$ \\
$C_2 \cong \mathfrak{sp}_4$
  & $10$ & $3$ & $5(k{+}3)/3$ & $20$ & L & $3$ \\
$G_2$
  & $14$ & $4$ & $7(k{+}4)/4$ & $28$ & L & $3$ \\
$F_4$
  & $52$ & $9$ & $26(k{+}9)/9$ & $104$ & L & $3$ \\
$E_6$
  & $78$ & $12$ & $13(k{+}12)/4$ & $156$ & L & $3$ \\
$E_7$
  & $133$ & $18$ & $133(k{+}18)/36$ & $266$ & L & $3$ \\
$E_8$
  & $248$ & $30$ & $62(k{+}30)/15$ & $496$ & L & $3$ \\
\bottomrule
\end{tabular}
}
\end{table}

\noindent
The kappa formula is universal:
$\kappa(V_k(\fg)) = \dim(\fg)\,(k + h^\vee)/(2h^\vee)$.
The central charge sum $c + c' = 2\dim(\fg)$ under the
Feigin--Frenkel involution $k \mapsto -k - 2h^\vee$, and
the complementarity $\kappa + \kappa' = 0$ holds for all
affine Kac--Moody algebras.  For non-simply-laced types,
the dual Coxeter number $h^\vee$ differs from the Coxeter
number $h$ ($h^\vee = 3$ for $B_2$ and $C_2$, while $h = 4$).
The shadow class is~L regardless: the Jacobi identity kills the
quartic on every primary line.

\begin{remark}[$B_2$/$C_2$ isomorphism]
The Lie algebras $B_2 = \mathfrak{so}_5$ and
$C_2 = \mathfrak{sp}_4$ are isomorphic, sharing $\dim = 10$,
$h^\vee = 3$, and $h = 4$.  Their Cartan matrices are transposes
of each other, but all shadow invariants coincide.
\end{remark}


\subsection{Complementarity table}\label{ssec:complementarity-table}

\begin{table}[ht]
\centering
\caption{Complementarity data: $\kappa(\cA) + \kappa(\cA^!)$ and
$c(\cA) + c(\cA^!)$ under Koszul duality}
\label{tab:complementarity}
\renewcommand{\arraystretch}{1.35}
{\small
\begin{tabular}{llccc}
\toprule
\textbf{Family $\cA$}
  & \textbf{Koszul dual $\cA^!$}
  & $\boldsymbol{c + c'}$
  & $\boldsymbol{\kappa + \kappa'}$
  & \textbf{Self-dual point} \\
\midrule
$\cH_k$
  & $\mathrm{Sym}^{\mathrm{ch}}(V^*)$
  & --- & $0$ & none \\
$V_\Lambda$ (lattice)
  & $V_{\Lambda^!}$
  & --- & $0$ & $\Lambda$ self-dual \\
Free fermion
  & $\mathrm{Sym}^{\mathrm{ch}}(\gamma)$
  & $0$ & $0$ & --- \\
$\beta\gamma_\lambda$
  & $bc_\lambda$
  & $0$ & $0$ & --- \\[3pt]
\midrule
$V_k(\mathfrak{sl}_2)$
  & $V_{-k-4}(\mathfrak{sl}_2)$
  & $6$ & $0$ & $k = -2$ (critical) \\
$V_k(\mathfrak{sl}_N)$
  & $V_{-k-2N}(\mathfrak{sl}_N)$
  & $2(N^2{-}1)$ & $0$ & $k = -N$ (critical) \\
$V_k(\fg)$ (general)
  & $V_{-k-2h^\vee}(\fg)$
  & $2\dim(\fg)$ & $0$ & $k = -h^\vee$ (critical) \\[3pt]
\midrule
$\mathrm{Vir}_c$
  & $\mathrm{Vir}_{26-c}$
  & $26$ & $13$ & $c = 13$ \\
$\cW_3$ at $c$
  & $\cW_3$ at $100{-}c$
  & $100$ & $250/3$ & $c = 50$ \\
$\cW_4$ at $c$
  & $\cW_4$ at $246{-}c$
  & $246$ & $533/2$ & $c = 123$ \\
$\cW_5$ at $c$
  & $\cW_5$ at $488{-}c$
  & $488$ & $9394/15$ & $c = 244$ \\
$\cW_N$ at $c$ (general)
  & $\cW_N$ at $K_N{-}c$
  & $K_N$ & $(H_N{-}1)\,K_N$ & $c = K_N/2$ \\
\bottomrule
\end{tabular}
}
\end{table}

\noindent
Here $K_N = 2(N{-}1) + 4(N^2{-}1)N$ is the Koszul conductor for
$\cW_N = \cW(\mathfrak{sl}_N)$ (the sum $c + c'$ under the
Feigin--Frenkel involution $k \mapsto -k - 2N$ on the underlying
$\mathfrak{sl}_N$).  The key structural distinction: for all
affine Kac--Moody algebras, $\kappa + \kappa' = 0$ (anti-symmetry
under the Feigin--Frenkel involution); for $\cW$-algebras,
$\kappa + \kappa' = (H_N - 1) \cdot K_N \neq 0$.
The Virasoro ($N = 2$) gives $\kappa + \kappa' = \tfrac{1}{2} \cdot 26
= 13$, which vanishes only at the self-dual point $c = 13$
where $\kappa(\mathrm{Vir}_{13}) = \kappa(\mathrm{Vir}_{13}) = 13/2$.


\subsection{Genus tables}\label{ssec:genus-tables}

The Faber--Pandharipande numbers and the genus expansion
$F_g(\cA) = \kappa(\cA) \cdot \lambda_g^{\mathrm{FP}}$.

\begin{table}[ht]
\centering
\caption{Faber--Pandharipande intersection numbers $\lambda_g^{\mathrm{FP}}$}
\label{tab:faber-pandharipande}
\renewcommand{\arraystretch}{1.35}
{\small
\begin{tabular}{ccc}
\toprule
$g$ & $\lambda_g^{\mathrm{FP}}$ & Decimal \\
\midrule
$1$ & $1/24$ & $0.041\overline{6}$ \\
$2$ & $7/5760$ & $1.215 \times 10^{-3}$ \\
$3$ & $31/967680$ & $3.204 \times 10^{-5}$ \\
$4$ & $127/154828800$ & $8.201 \times 10^{-7}$ \\
$5$ & $73/3503554560$ & $2.084 \times 10^{-8}$ \\
\bottomrule
\end{tabular}
}
\end{table}

\noindent
The formula is
$\lambda_g^{\mathrm{FP}} = \frac{2^{2g-1} - 1}{2^{2g-1}}
\cdot \frac{|B_{2g}|}{(2g)!}$,
where $B_{2g}$ is the $2g$-th Bernoulli number.  The generating
function is $\sum_{g \geq 1} \lambda_g^{\mathrm{FP}}\, \hbar^{2g}
= \hat{A}(i\hbar) - 1$, where $\hat{A}(x) = (x/2)/\sin(x/2)$ is
the $\hat{A}$-genus.  The numbers $\lambda_g^{\mathrm{FP}}$ are
strictly positive, strictly decreasing, and decay like
$(2\pi)^{-2g}$.

\begin{table}[ht]
\centering
\caption{Genus expansion $F_g(\cA) = \kappa \cdot \lambda_g^{\mathrm{FP}}$
for selected families}
\label{tab:genus-expansion}
\renewcommand{\arraystretch}{1.35}
{\small
\begin{tabular}{lccccc}
\toprule
\textbf{Algebra $\cA$}
  & $\boldsymbol{\kappa}$
  & $\boldsymbol{F_1}$
  & $\boldsymbol{F_2}$
  & $\boldsymbol{F_3}$ \\
\midrule
$\cH_1$ (Heisenberg, $k{=}1$)
  & $1$ & $1/24$ & $7/5760$ & $31/967680$ \\
$V_1(\mathfrak{sl}_2)$
  & $9/4$ & $3/32$ & $7/2560$ & $31/430080$ \\
$V_1(\mathfrak{sl}_3)$
  & $16/3$ & $2/9$ & $7/1080$ & $31/181440$ \\
$V_\Lambda$ (Leech, $\operatorname{rank}{=}24$)
  & $24$ & $1$ & $7/240$ & $31/40320$ \\
$\beta\gamma$ (standard, $\lambda{=}1$)
  & $1$ & $1/24$ & $7/5760$ & $31/967680$ \\
$\mathrm{Vir}_1$
  & $1/2$ & $1/48$ & $7/11520$ & $31/1935360$ \\
$\mathrm{Vir}_{13}$
  & $13/2$ & $13/48$ & $91/11520$ & $403/1935360$ \\
$\mathrm{Vir}_{26}$
  & $13$ & $13/24$ & $91/5760$ & $403/967680$ \\
$\cW_3$ at $c{=}6$
  & $5$ & $5/24$ & $7/1152$ & $31/193536$ \\
\bottomrule
\end{tabular}
}
\end{table}

\noindent
The entries for $V_1(\mathfrak{sl}_3)$ use
$\kappa = (8 \cdot 4)/(2 \cdot 3) = 16/3$.  The Leech lattice
($\operatorname{rank} = 24$) gives $F_1 = 1$, reflecting the
Jacobi $\Delta$-function origin of the Leech theta series.

\subsection{Shadow growth rate table}\label{ssec:growth-rate-table}

\begin{table}[ht]
\centering
\caption{Shadow growth rate $\rho(\cA)$ for class~M algebras
(Virasoro at selected central charges)}
\label{tab:growth-rate}
\renewcommand{\arraystretch}{1.35}
{\small
\begin{tabular}{ccccl}
\toprule
$c$
  & $\kappa$
  & $Q^{\mathrm{contact}}$
  & $\rho$
  & \textbf{Regime} \\
\midrule
$1/2$ (Ising)
  & $1/4$ & $40/49$ & $\approx 12.53$ & divergent \\
$1$ (free boson)
  & $1/2$ & $10/27$ & $\approx 6.24$ & divergent \\
$4$
  & $2$ & $5/84$ & $\approx 1.54$ & divergent \\
$c^* \approx 6.125$
  & $\approx 3.06$ & $\approx 0.031$ & $1$ & critical \\
$13$ (self-dual)
  & $13/2$ & $10/1131$ & $\approx 0.467$ & convergent \\
$25$
  & $25/2$ & $2/735$ & $\approx 0.242$ & convergent \\
$26$ (bosonic string)
  & $13$ & $5/1976$ & $\approx 0.232$ & convergent \\
\bottomrule
\end{tabular}
}
\end{table}

\noindent
The shadow growth rate $\rho = \sqrt{(9\alpha^2 + 2\Delta)/(4\kappa^2)}$
determines the convergence of the shadow tower.  For Virasoro,
$\rho^2 = (180c + 872)/[c^2(5c+22)]$.  The critical central charge
$c^* \approx 6.125$ (the positive real root of~\eqref{eq:critical-cubic})
separates the convergent regime ($c > c^*$, tower summable) from the
divergent regime ($c < c^*$, tower grows exponentially).  Class~G and
class~L algebras have $\rho = 0$ (the tower terminates, so there is no
exponential growth).

\subsection{Shadow field table}\label{ssec:shadow-field-table}

For class~$\mathbf{M}$ algebras at rational~$c$, the shadow metric
defines an imaginary quadratic field $K_L = \mathbb{Q}(\sqrt{D_0})$
through which the Epstein zeta factors.  The discriminant
$\operatorname{disc}(Q_L) = -32\kappa^2\Delta$ is negative for
positive~$c$.

\begin{table}[ht]
\centering
\caption{Shadow fields for the Virasoro algebra at rational~$c$.
The squarefree kernel $D_0$ is computed from
$\operatorname{disc}(Q_L) = -320c^2/(5c+22)$;
$K_L = \mathbb{Q}(\sqrt{D_0})$.}
\label{tab:shadow-fields-vir}
\renewcommand{\arraystretch}{1.35}
{\small
\begin{tabular}{lcccccl}
\toprule
\textbf{Model} & $c$ & $\kappa$ & $\Delta$
  & $D_0$
  & $h(D_0)$
  & \textbf{Shadow field} \\
\midrule
Ising $\mathcal{M}(3,4)$ & $1/2$ & $1/4$ & $80/49$
  & $-10$ & $2$
  & $\mathbb{Q}(\sqrt{-10})$ \\
Tricritical Ising & $7/10$ & $7/20$ & $200/309$
  & $-510$ & ---
  & $\mathbb{Q}(\sqrt{-510})$ \\
Free boson & $1$ & $1/2$ & $40/27$
  & $-15$ & $2$
  & $\mathbb{Q}(\sqrt{-15})$ \\
$c = 4$ & $4$ & $2$ & $20/21$
  & $-210$ & ---
  & $\mathbb{Q}(\sqrt{-210})$ \\
$c^* \approx 6.12$ & $c^*$ & $\approx 3.06$ & $\approx 0.99$
  & ---  & ---
  & (critical) \\
Self-dual & $13$ & $13/2$ & $40/87$
  & $-435$ & ---
  & $\mathbb{Q}(\sqrt{-435})$ \\
Critical string & $26$ & $13$ & $40/152$
  & $-190$ & ---
  & $\mathbb{Q}(\sqrt{-190})$ \\
\bottomrule
\end{tabular}
}
\end{table}

\noindent
Here $D_0$ is the squarefree kernel of the discriminant,
$h(D_0)$ is the class number of~$\mathbb{Q}(\sqrt{D_0})$,
and the discriminant formula is
$\operatorname{disc}(Q_L) = -320\,c^2/(5c+22)$, so
$D_0 = \operatorname{sqfree}(-5\,q\,(5p + 22q))$ for
$c = p/q$ in lowest terms.  The Ising model ($D_0 = -10$)
and the free boson ($D_0 = -15$) both have class number~$2$;
by Davenport--Heilbronn, the corresponding Epstein zeta
functions may have zeros off the critical line.


\subsection{Virasoro shadow tower table}\label{ssec:virasoro-tower-table}

\begin{table}[ht]
\centering
\caption{Virasoro shadow coefficients $S_r$ for $r = 2, \dotsc, 10$}
\label{tab:virasoro-tower}
\renewcommand{\arraystretch}{1.35}
{\small
\begin{tabular}{cl}
\toprule
$r$ & $S_r(\mathrm{Vir}_c)$ \\
\midrule
$2$ & $c/2$ \\
$3$ & $2$ \\
$4$ & $10/[c(5c+22)]$ \\
$5$ & $-48/[c^2(5c+22)]$ \\
$6$ & $80(45c+193)/[3c^3(5c+22)^2]$ \\
$7$ & $-2880(15c+61)/[7c^4(5c+22)^2]$ \\
$8$ & $80(2025c^2+16470c+33314)/[c^5(5c+22)^3]$ \\
$9$ & $-1280(2025c^2+15570c+29554)/[3c^6(5c+22)^3]$ \\
$10$ & $256(91125c^3+1050975c^2+3989790c+4969967)/[c^7(5c+22)^4]$ \\
\bottomrule
\end{tabular}
}
\end{table}

\noindent
Each $S_r$ is a rational function of $c$ with poles at $c = 0$ and
$c = -22/5$; these are the two values where the quadratic Casimir
$c(5c+22)$ vanishes.  The sign pattern
$S_r > 0$ for even $r$ and $S_r < 0$ for odd $r \geq 5$
(at $c > 0$) reflects the oscillatory decay $\cos(r\theta + \phi)$
from the complex branch points of the shadow generating function.
All entries have been verified by two independent methods: closed-form
derivation and the convolution recursion from $H^2 = t^4 Q_L$.


%======================================================================
%  COMPUTATIONS.TEX -- Sections 7--8 of the standalone paper
%  The genus-2 free energy of W_3 and the gravitational coproduct
%======================================================================


%======================================================================
\section{The genus-2 free energy of \texorpdfstring{$W_3$}{W3}}
\label{sec:genus2-w3}
%======================================================================

The preceding sections developed the theory in the single-channel
regime: one strong generator, one primary line, one modular
characteristic~$\kappa$.  We now compute the first explicit
multi-channel invariant.  The $W_3$ algebra has two strong
generators---the stress tensor~$T$ of conformal weight~$2$ and the
spin-3 current~$W$ of conformal weight~$3$---and the genus-2 free
energy $F_2(W_3)$ receives contributions from \emph{both} channels
and their interactions.

\subsection{The Frobenius algebra of the bar-complex CohFT}
\label{subsec:w3-frobenius}

The bar-complex CohFT of $W_3$ has a two-dimensional underlying
Frobenius algebra with basis $\{e_T, e_W\}$.  The metric is diagonal,
determined by the per-channel modular characteristics:
\begin{equation}\label{eq:w3-metric}
\eta_{TT} = \kappa_T = \frac{c}{2},
\qquad
\eta_{WW} = \kappa_W = \frac{c}{3},
\qquad
\eta_{TW} = 0.
\end{equation}
The total modular characteristic is $\kappa(W_3) = \kappa_T + \kappa_W
= 5c/6 = c \cdot (H_3 - 1)$, where $H_3 = 1 + 1/2 + 1/3 = 11/6$
is the third harmonic number.  The inverse metric gives the
propagators:
\begin{equation}\label{eq:w3-propagators}
\eta^{TT} = \frac{2}{c},
\qquad
\eta^{WW} = \frac{3}{c}.
\end{equation}
Both propagators have weight~$1$:
the bar-complex propagator is $d\!\log E(z,w)$, where
$E(z,w)$ is the prime form (a section of
$K^{-1/2} \boxtimes K^{-1/2}$), so $d\!\log E$ has weight~$1$
regardless of the conformal weight of the field being sewed.

The three-point structure constants are determined by the OPE.
The $\mathbb{Z}_2$ symmetry $W \mapsto -W$ forces all three-point
functions with an odd number of $W$-insertions to vanish:
\begin{equation}\label{eq:w3-C3}
C_{TTT} = C_{TWW} = C_{WWT} = c,
\qquad
C_{TTW} = C_{TWT} = C_{WTT} = C_{WWW} = 0.
\end{equation}
The derivations are:
$T_{(1)}T = 2T$ gives $C^T_{TT} = 2$, whence
$C_{TTT} = \eta_{TT} \cdot 2 = c$;
$T_{(1)}W = 3W$ gives $C^W_{TW} = 3$, whence
$C_{TWW} = \eta_{WW} \cdot 3 = c$;
$W_{(3)}W = 2T$ gives $C^T_{WW} = 2$, whence
$C_{WWT} = \eta_{TT} \cdot 2 = c$.

The genus-0 four-point vertex exhibits a remarkable universality.
For any channel pair $(i,j) \in \{T,W\}^2$:
\begin{equation}\label{eq:w3-V04}
V_{0,4}(i,i \mid j,j)
\;=\;
\sum_m \eta^{mm}\, C_{iim}\, C_{jjm}
\;=\;
\eta^{TT} \cdot c \cdot c
\;=\;
2c.
\end{equation}
Only the $T$-channel contributes to the sum: $C_{iiT} = c$ for
$i \in \{T,W\}$, while $C_{iiW} = 0$ for both $i = T$ (since
$C_{TTW} = 0$) and $i = W$ (since $C_{WWW} = 0$).
The $W$-channel is invisible at genus-0 valence~4.


\subsection{The seven stable graphs at genus 2}
\label{subsec:genus2-graphs}

The moduli space $\overline{\mathcal{M}}_{2,0}$ has a
stratification by seven stable graphs $\Gamma_0, \ldots, \Gamma_6$:

\begin{center}
\small
\renewcommand{\arraystretch}{1.3}
\begin{tabular}{clllccl}
& \textbf{Name}
& \textbf{Vertices}
& \textbf{Edges}
& $|\mathrm{Aut}|$
& $h^1$
& \textbf{Type} \\
\hline
$\Gamma_0$ & smooth
  & $(g{=}2,\, \mathrm{val}{=}0)$
  & $0$
  & $1$ & $0$ & compact stratum \\[2pt]
$\Gamma_1$ & figure-eight
  & $(g{=}1,\, \mathrm{val}{=}2)$
  & $1$ self-loop
  & $2$ & $1$ & one-loop \\[2pt]
$\Gamma_2$ & banana
  & $(g{=}0,\, \mathrm{val}{=}4)$
  & $2$ self-loops
  & $8$ & $2$ & two-loop sunset \\[2pt]
$\Gamma_3$ & dumbbell
  & $(g{=}1,\, \mathrm{val}{=}1) + (g{=}1,\, \mathrm{val}{=}1)$
  & $1$ bridge
  & $2$ & $0$ & separating \\[2pt]
$\Gamma_4$ & theta
  & $(g{=}0,\, \mathrm{val}{=}3) + (g{=}0,\, \mathrm{val}{=}3)$
  & $3$ bridges
  & $12$ & $2$ & non-separating \\[2pt]
$\Gamma_5$ & lollipop
  & $(g{=}0,\, \mathrm{val}{=}3) + (g{=}1,\, \mathrm{val}{=}1)$
  & $1$ self-loop $+$ $1$ bridge
  & $2$ & $1$ & mixed
\end{tabular}
\end{center}
Every vertex satisfies the stability condition $2g_v + \mathrm{val}_v
\geq 3$, and the genus formula $h^1(\Gamma) + \sum_v g_v = 2$ holds
in each case.  Note that the smooth graph $\Gamma_0$ has no edges and
contributes the integrated class on $\mathcal{M}_{2,0}$ itself; the
remaining five graphs are the boundary strata.


\subsection{Multi-channel Feynman rules}
\label{subsec:w3-feynman-rules}

For each graph $\Gamma$ with $|E(\Gamma)|$ edges, a \emph{channel
assignment} is a function $\sigma \colon E(\Gamma) \to \{T, W\}$.
The Feynman rules of the bar-complex CohFT assign to each pair
$(\Gamma, \sigma)$ an amplitude
\begin{equation}\label{eq:feynman-rule}
\mathcal{A}(\Gamma, \sigma)
\;=\;
\frac{1}{|\mathrm{Aut}(\Gamma)|}
\prod_{e \in E(\Gamma)} \eta^{\sigma(e),\, \sigma(e)}
\prod_{v \in V(\Gamma)} V_v\bigl(g_v;\, \sigma|_{\mathrm{half\text{-}edges}(v)}\bigr),
\end{equation}
where the vertex factors are:
\begin{itemize}
\item \emph{Genus-0, trivalent}: $V_v = C_{\sigma(e_1)\sigma(e_2)\sigma(e_3)}$.
\item \emph{Genus-0, four-valent} (the banana vertex):
  $V_v = \sum_m \eta^{mm} C_{\sigma(e_1)\sigma(e_2)m}\,
         C_{m\,\sigma(e_3)\sigma(e_4)}$.
\item \emph{Genus-1, valence $n$}:
  $V_v = \kappa_{\sigma(e)}/24$ (per-channel genus-1 universality,
  proved unconditionally).
\end{itemize}
Each half-edge at a self-loop contributes a pair of half-edges to
the same vertex, both carrying channel~$\sigma(e)$.


\subsection{Per-graph amplitudes}
\label{subsec:per-graph-amplitudes}

We now compute the amplitude of each boundary graph, decomposed into
\emph{diagonal} (all edges carry the same channel) and \emph{mixed}
(at least two edges carry different channels) contributions.

\subsubsection*{$\Gamma_1$: figure-eight}

One edge, two channel assignments.  The vertex is genus-1 with
valence~2, so the vertex factor for channel~$i$ is $\kappa_i/24$:
\begin{align}
\mathcal{A}(\Gamma_1, T) &= \tfrac{1}{2}\, \eta^{TT} \cdot
  \tfrac{\kappa_T}{24}
  = \tfrac{1}{2} \cdot \tfrac{2}{c} \cdot \tfrac{c}{48}
  = \tfrac{1}{48}, \label{eq:fig-eight-T}\\
\mathcal{A}(\Gamma_1, W) &= \tfrac{1}{2}\, \eta^{WW} \cdot
  \tfrac{\kappa_W}{24}
  = \tfrac{1}{2} \cdot \tfrac{3}{c} \cdot \tfrac{c}{72}
  = \tfrac{1}{48}. \label{eq:fig-eight-W}
\end{align}
Both channels contribute equally: the propagator~$1/\kappa_i$ cancels
the genus-1 vertex factor~$\kappa_i/24$, leaving the universal value
$1/48$ independent of $c$.  The total is $1/24 = \lambda_1^{\mathrm{FP}}$.
No mixed channel assignments exist (single edge).

\subsubsection*{$\Gamma_2$: banana}

Two edges, four channel assignments, $|\mathrm{Aut}| = 8$.  The
vertex is genus-0 with valence~4, and the four-point vertex factor
is $V_{0,4}(i,i \mid j,j) = 2c$ universally~\eqref{eq:w3-V04}.
\begin{align}
\mathcal{A}(\Gamma_2, TT) &=
  \tfrac{1}{8}\, (\eta^{TT})^2 \cdot 2c
  = \tfrac{1}{8} \cdot \tfrac{4}{c^2} \cdot 2c
  = \frac{1}{c}, \label{eq:banana-TT}\\
\mathcal{A}(\Gamma_2, WW) &=
  \tfrac{1}{8}\, (\eta^{WW})^2 \cdot 2c
  = \tfrac{1}{8} \cdot \tfrac{9}{c^2} \cdot 2c
  = \frac{9}{4c}, \label{eq:banana-WW}\\
\mathcal{A}(\Gamma_2, TW{+}WT) &=
  \tfrac{1}{8} \cdot 2 \cdot \eta^{TT}\, \eta^{WW} \cdot 2c
  = \tfrac{1}{8} \cdot 2 \cdot \tfrac{6}{c^2} \cdot 2c
  = \frac{3}{c}. \label{eq:banana-mixed}
\end{align}
The cross-channel contribution is $\delta_{\mathrm{banana}} = 3/c$.

\subsubsection*{$\Gamma_3$: dumbbell}

One bridge edge, two channel assignments, $|\mathrm{Aut}| = 2$.
Each vertex has genus~1 and valence~1, with vertex factor
$\kappa_i/24$:
\begin{align}
\mathcal{A}(\Gamma_3, T) &=
  \tfrac{1}{2}\, \eta^{TT} \cdot
  \bigl(\tfrac{\kappa_T}{24}\bigr)^2
  = \tfrac{1}{2} \cdot \tfrac{2}{c} \cdot \tfrac{c^2}{2304}
  = \frac{c}{2304}, \label{eq:dumbbell-T}\\
\mathcal{A}(\Gamma_3, W) &=
  \tfrac{1}{2}\, \eta^{WW} \cdot
  \bigl(\tfrac{\kappa_W}{24}\bigr)^2
  = \tfrac{1}{2} \cdot \tfrac{3}{c} \cdot \tfrac{c^2}{5184}
  = \frac{c}{3456}. \label{eq:dumbbell-W}
\end{align}
The total is $c/2304 + c/3456 = (3c + 2c)/6912 = 5c/6912 =
\kappa(W_3)/1152$.  No mixed channel assignments exist (single edge).

\subsubsection*{$\Gamma_4$: theta}

Three bridge edges, eight channel assignments, $|\mathrm{Aut}| = 12$.
Each vertex has genus~0 and valence~3, with vertex factor
$C_{\sigma(e_1)\sigma(e_2)\sigma(e_3)}$.
\begin{itemize}
\item \emph{All-$T$}: both vertices get $C_{TTT} = c$.
  Amplitude $= \frac{1}{12}\, (\eta^{TT})^3 \cdot c^2
  = \frac{1}{12} \cdot \frac{8}{c^3} \cdot c^2
  = \frac{2}{3c}$.

\item \emph{All-$W$}: both vertices get $C_{WWW} = 0$.
  Amplitude $= 0$.

\item \emph{Type $(T,T,W)$} (three such assignments):
  each vertex receives one~$W$ half-edge, giving $C_{TTW} = 0$.
  All three vanish.

\item \emph{Type $(T,W,W)$} (three such assignments):
  each vertex receives half-edges $(T,W,W)$, giving
  $C_{TWW} = c$.  The propagator product is
  $\eta^{TT} (\eta^{WW})^2 = (2/c)(9/c^2) = 18/c^3$.
  Each assignment contributes $18c^2/c^3 = 18/c$; with
  $|\mathrm{Aut}| = 12$, the total mixed amplitude is
  \begin{equation}\label{eq:theta-mixed}
  \delta_{\mathrm{theta}} = \frac{3 \cdot 18}{12c}
  = \frac{9}{2c}.
  \end{equation}
\end{itemize}

\subsubsection*{$\Gamma_5$: lollipop}

One self-loop at vertex~$v_0$ (genus~0, valence~3) and one bridge
to vertex~$v_1$ (genus~1, valence~1).  Two edges, four channel
assignments, $|\mathrm{Aut}| = 2$.

\begin{itemize}
\item \emph{$(T,T)$}: self-loop~$T$, bridge~$T$.  Vertex~$v_0$ gets
  half-edges $(T,T,T) \to C_{TTT} = c$.  Vertex~$v_1$ gets
  half-edge $T \to \kappa_T/24 = c/48$.  Propagators:
  $(\eta^{TT})^2 = 4/c^2$.  With $|\mathrm{Aut}|$:
  $\frac{1}{2} \cdot \frac{4}{c^2} \cdot c \cdot \frac{c}{48}
  = \frac{1}{24}$.

\item \emph{$(W,W)$}: self-loop~$W$, bridge~$W$.  Vertex~$v_0$ gets
  $(W,W,W) \to C_{WWW} = 0$.  Vanishes.

\item \emph{$(T,W)$}: self-loop~$T$, bridge~$W$.  Vertex~$v_0$ gets
  $(T,T,W) \to C_{TTW} = 0$.  Vanishes.

\item \emph{$(W,T)$}: self-loop~$W$, bridge~$T$.  Vertex~$v_0$ gets
  $(W,W,T) \to C_{WWT} = c$.  Vertex~$v_1$ gets
  $T \to \kappa_T/24 = c/48$.  Propagators:
  $\eta^{WW} \cdot \eta^{TT} = (3/c)(2/c) = 6/c^2$.
  With $|\mathrm{Aut}|$:
  \begin{equation}\label{eq:lollipop-mixed}
  \delta_{\mathrm{lollipop}}
  = \tfrac{1}{2} \cdot \tfrac{6}{c^2} \cdot c \cdot \tfrac{c}{48}
  = \tfrac{6c^2}{96c^2} = \frac{1}{16}.
  \end{equation}
\end{itemize}
The lollipop mixed amplitude $1/16$ is \emph{independent of~$c$}:
the $c^2$ in the numerator (from $C_{WWT} \cdot \kappa_T/24$)
cancels the $c^2$ in the denominator (from $\eta^{WW} \cdot \eta^{TT}$).


\subsection{The cross-channel correction}
\label{subsec:cross-channel}

Collecting the mixed-channel amplitudes from all boundary graphs:
\begin{equation}\label{eq:cross-channel-total}
\delta F_2
\;=\;
\underbrace{\frac{3}{c}}_{\text{banana}}
\;+\;
\underbrace{\frac{9}{2c}}_{\text{theta}}
\;+\;
\underbrace{\frac{1}{16}}_{\text{lollipop}}
\;+\;
\underbrace{\frac{21}{4c}}_{\text{barbell}}
\;=\;
\frac{51}{4c} + \frac{1}{16}
\;=\;
\frac{c + 204}{16c}.
\end{equation}
The figure-eight and dumbbell graphs contribute zero mixed amplitude
(single edges admit only diagonal assignments).  The correction
decomposes into two structurally distinct pieces:
\begin{itemize}
\item A \emph{$c$-dependent} piece $51/(4c)$, arising from
  genus-0 vertex factors on the banana ($3/c$), theta ($9/(2c)$),
  and barbell ($21/(4c)$) graphs.  This is
  $R$-matrix independent: genus-0 vertices receive no corrections
  from the CohFT $R$-matrix.
\item A \emph{$c$-independent} piece $1/16$, arising from the
  lollipop graph, which mixes a genus-0 vertex factor with a
  genus-1 vertex factor $\kappa_i/24$.
\end{itemize}

\begin{remark}[Comparison with the scalar approximation]
\label{rem:scalar-comparison}
The per-channel (diagonal) graph sum, computed by applying Theorem~D
to each channel independently, gives
\begin{equation}\label{eq:F2-scalar}
F_2^{\mathrm{scal}}(W_3) = \kappa \cdot \lambda_2^{\mathrm{FP}}
= \frac{5c}{6} \cdot \frac{7}{5760} = \frac{7c}{6912}.
\end{equation}
The ratio of the naive cross-channel correction to the scalar
approximation is
\begin{equation}\label{eq:cross-channel-ratio}
\frac{\delta F_2}{F_2^{\mathrm{scal}}}
= \frac{432(c+204)}{7c^2},
\end{equation}
which diverges as $c \to 0$ and decays as $c^{-1}$ at large~$c$.
At $c = 4$: the ratio is approximately~$802$, and the cross-channel
correction \emph{dominates} the scalar term.  At $c = 50$: the
ratio is approximately~$6.3$.
At large~$c$, the lollipop's $c$-independent contribution $1/16$
dominates, and the correction becomes a subleading $O(1/c)$ effect
relative to the $O(c)$ scalar term.
\end{remark}

\begin{remark}[Status of the cross-channel correction]
\label{rem:cross-channel-status}
The computation above uses the naive genus-$1$ vertex factors
$\kappa_i/24$ without $R$-matrix corrections.  The
multi-generator universality problem
(\textbf{op:multi-generator-universality}) is resolved negatively:
the CohFT $R$-matrix does \emph{not} produce compensating terms
that cancel $\delta F_2$.  The cross-channel correction
$\delta F_2 = (c{+}204)/(16c)$ is a genuine additional term,
and the correct decomposition is
$F_g = \kappa \cdot \lambda_g^{\mathrm{FP}} +
\delta F_g^{\mathrm{cross}}$.
\end{remark}


\subsection{The complete amplitude table}
\label{subsec:amplitude-table}

We record the full per-graph data:

\begin{center}
\small
\renewcommand{\arraystretch}{1.3}
\begin{tabular}{lcccc}
\textbf{Graph}
& \textbf{Diagonal ($T$)}
& \textbf{Diagonal ($W$)}
& \textbf{Mixed}
& \textbf{Total} \\
\hline
$\Gamma_1$ figure-eight
& $1/48$ & $1/48$ & $0$ & $1/24$ \\[2pt]
$\Gamma_2$ banana
& $1/c$ & $9/(4c)$ & $3/c$ & $25/(4c)$ \\[2pt]
$\Gamma_3$ dumbbell
& $c/2304$ & $c/3456$ & $0$ & $5c/6912$ \\[2pt]
$\Gamma_4$ theta
& $2/(3c)$ & $0$ & $9/(2c)$ & $31/(6c)$ \\[2pt]
$\Gamma_5$ lollipop
& $1/24$ & $0$ & $1/16$ & $5/48$ \\
\end{tabular}
\end{center}

The table records the boundary-graph amplitudes only; the
smooth graph $\Gamma_0$ contributes an additional integrated
class on $\mathcal{M}_{2,0}$ whose per-channel decomposition
is determined by the constraint
$\sum_\Gamma \mathcal{A}(\Gamma) = \kappa \cdot \lambda_2^{\mathrm{FP}}$
when the full universality identity holds.
The mixed column sums to
$\delta F_2 = (c+204)/(16c)$.

Every entry in this table has been verified against the compute
module \texttt{w3\_genus2.py}, using exact rational arithmetic
(\texttt{fractions.Fraction}), at the seven test central charges
$c \in \{1, 2, 4, 13, 26, 50, 100\}$.


\subsection{Planted-forest corrections}
\label{subsec:planted-forest}

The planted-forest correction on each primary line is
$\delta_{\mathrm{pf}}^{(2,0)} = S_3(10 S_3 - \kappa)/48$.

\emph{$T$-line}: $S_3^T = 2$ (the Virasoro cubic shadow),
$\kappa_T = c/2$.
\begin{equation}\label{eq:pf-T}
\delta_{\mathrm{pf}}^T
= \frac{2(20 - c/2)}{48}
= \frac{40 - c}{48}.
\end{equation}
This vanishes at $c = 40$, is positive for $c < 40$, and negative
for $c > 40$.

\emph{$W$-line}: $S_3^W = 0$ by the $\mathbb{Z}_2$ parity
$W \mapsto -W$, which forces all odd-arity shadows on the $W$-line
to vanish.  Hence $\delta_{\mathrm{pf}}^W = 0$.

The total planted-forest correction is $\delta_{\mathrm{pf}} =
(40-c)/48$, living entirely on the $T$-line.


\subsection{Koszul duality constraints}
\label{subsec:w3-koszul-duality}

Under the Koszul involution, $W_3$ at central charge~$c$ is
dual to $W_3$ at central charge $c' = 100 - c$
(the Feigin--Frenkel involution for $\mathfrak{sl}_3$).  The modular
characteristics satisfy
\begin{equation}
\kappa(c) + \kappa(100-c) = \frac{5c}{6} + \frac{5(100-c)}{6}
= \frac{250}{3},
\end{equation}
and the scalar free energies satisfy $F_2^{\mathrm{scal}}(c) +
F_2^{\mathrm{scal}}(100-c) = (250/3) \cdot \lambda_2^{\mathrm{FP}}$.
The cross-channel correction at the dual pair is
\[
\delta F_2(c) + \delta F_2(100-c)
= \frac{c+204}{16c} + \frac{304-c}{16(100-c)}
= -\,\frac{c^2 - 100c - 10200}{8c(100-c)}.
\]
This sum is not constant in~$c$: the cross-channel correction does
not respect Koszul additivity.  If universality holds ($\delta F_2$
is cancelled by $R$-matrix corrections), this is automatic.  If
the naive correction persists, it produces a Koszul-asymmetric
contribution to $F_2$---an observable departure from single-channel
behaviour.


%======================================================================
\section{The gravitational coproduct is primitive}
\label{sec:gravitational-coproduct}
%======================================================================

In gauge theory, the gluon is a fundamental field: it can split.
A single gluon entering a scattering region can exit as two gluons,
and the coproduct of the dg-shifted Yangian encodes this splitting
at the algebraic level.  In gravity, the graviton is not fundamental
but composite.  The stress tensor is the Sugawara quadratic
\begin{equation}\label{eq:sugawara}
T = \frac{1}{k+2}\Bigl(\tfrac{1}{2} {:}J^0 J^0{:}
+ {:}J^+ J^-{:}\Bigr),
\end{equation}
built from two copies of the affine connection by normal ordering.
The compositeness of the graviton has a precise algebraic
consequence: the transferred coproduct is trivial.


\subsection{The DS-HPL transfer theorem}
\label{subsec:ds-hpl}

The Drinfeld--Sokolov reduction $\hat{\mathfrak{sl}}_2 \to
\mathrm{Vir}$ is a BRST reduction with a strong deformation retract
$(i, p, h)$ from the BRST complex to the Virasoro cohomology.
The homotopy $h$ has ghost-number degree~$-1$: it maps
$V \mapsto b$ (ghost number $0 \to -1$) and
$c \mapsto U/2$ (ghost number $1 \to 0$).

\begin{theorem}[DS-HPL transfer]\label{thm:ds-hpl-standalone}
Let $\fg$ be a simple Lie algebra, and let
$Y^{\mathrm{dg}}(\widehat{\fg})$ denote the dg-shifted
Yangian at level~$k$, with $A_\infty$ products
$\{m_n^{\mathrm{aff}}\}$, coproduct
$\Delta_z^{\mathrm{aff}}$, and $r$-matrix
$r^{\mathrm{aff}}(z) = \Omega/z$.  The HPL transfer through the
principal BRST SDR $(i, p, h)$ from the DS complex to
$H^*_{\mathrm{BRST}} = \cW(\fg)$ produces:
\begin{enumerate}[label=\textup{(\roman*)}]
\item Transferred $A_\infty$ products
  $\{m_n^{\cW}\}_{n \geq 2}$ on the $\cW$-algebra sector,
  with $m_n^{\cW} \neq 0$ for all $n \geq 2$
  \textup{(}the entire infinite $A_\infty$ tower of class~$\mathbf{M}$%
  \textup{)}.
\item A transferred coproduct
  $\Delta_z^{\cW} \colon H_{\cW} \to
  H_{\cW} \otimes H_{\cW}$.
\item A transferred $r$-matrix
  $r^{\cW}(z)$ with leading pole of order
  $2h_{\max} - 1$, where $h_{\max}$ is the highest conformal
  weight among the generators of~$\cW(\fg)$.
  For $\fg = \mathfrak{sl}_2$:
  $r^{\mathrm{Vir}}(z) = (c/2)/z^3 + 2T/z$.
\end{enumerate}
\end{theorem}

The $r$-matrix deserves comment.  The affine $r$-matrix
$r^{\mathrm{aff}}(z) = \Omega/z$ is a simple pole---the
collision residue of the affine bar complex, whose OPE has at most
a double pole.  The Virasoro OPE has a quartic pole
($z^{-4}$ in $T(z)T(w)$); the bar kernel $d\!\log(z-w)$ absorbs
one power, producing the cubic pole $z^{-3}$ of
$r^{\mathrm{Vir}}(z)$.  This is the first instance of the general
principle: the $r$-matrix pole order equals the OPE pole order
minus one.

The key result is the following.

\begin{theorem}[Gravitational coproduct primitivity]
\label{thm:grav-primitivity-standalone}
For any simple Lie algebra~$\fg$, the HPL-transferred coproduct
of the principal $\cW$-algebra $\cW(\fg)$ is strictly primitive
at all arities:
\begin{equation}\label{eq:grav-prim}
\Delta_{z,k}^{\cW(\fg)} = 0
\qquad \text{for all } k \geq 2.
\end{equation}
In particular, for $\fg = \mathfrak{sl}_2$:
$\Delta_z^{\mathrm{Vir}}(T) = \tau_z(T) \otimes 1
+ 1 \otimes T$.  All nontrivial structure of the $\cW$-algebra
dg-shifted Yangian resides in the $r(z)$-twisted target product,
never in the coproduct.
\end{theorem}

\begin{proof}
The proof is a ghost-number obstruction.  The argument is
uniform in~$\fg$: only the ghost-number grading of the
principal BRST complex is used, and this grading is universal
for all principal DS reductions.  We write the proof for
general~$\fg$, illustrating with
$\fg = \mathfrak{sl}_2$ where explicit.

\emph{Step~1 (linear primitivity).}
The linear transferred coproduct
$(p \otimes p)(\Delta_z(\phi))$ is already primitive for
every generator~$\phi$ of $\cW(\fg)$.  The affine coproduct
$\Delta_z(\phi)$ involves the Casimir element
$\Omega = \sum_a J^a \otimes J_a$, and the projection~$p$
onto the $\cW$-algebra cohomology $H^*_{\mathrm{BRST}}$
annihilates the affine currents~$J^a$ (which have conformal
weight~$1$ and ghost number~$0$, but lie outside the
$\cW$-algebra sector).  Since every summand of~$\Omega$
contains at least one $J^a$ factor, $(p \otimes p)(\Omega) = 0$.

For $\fg = \mathfrak{sl}_2$: $p$ annihilates $J^0$, $J^\pm$;
the Casimir $\Omega(w{+}z,w) = J^0 \otimes J^0
+ J^+ \otimes J^- + J^- \otimes J^+$ is killed by
$p \otimes p$.

\emph{Step~2 (ghost-number obstruction at higher arity).}
At arity $k \geq 2$, every HPL tree has at least one internal
edge labelled by the homotopy~$h$.  The ghost-number accounting:
\begin{center}
\small
\begin{tabular}{lc}
Operation & Ghost degree \\
\hline
$h$ (homotopy) & $-1$ \\
$\delta$ (BRST perturbation) & $+1$ \\
$\mu$ (vertex algebra product) & $0$ \\
$\Delta_z$ (coproduct) & $0$ \\
\end{tabular}
\end{center}
The product $h\delta$ has ghost degree $-1 + 1 = 0$, but the
tree topology at arity $k \geq 2$ always contains one
\emph{uncompensated}~$h$: the outermost homotopy that is not
preceded by a $\delta$ insertion.  The overall output before
projection has ghost number $\leq -1$.  This counting is
independent of~$\fg$: it depends only on the tree
structure of the HPL formula and the universal ghost-number
assignments.

\emph{Step~3 (projection kills everything).}
The $\cW$-algebra cohomology sits in ghost number~$0$.
The projection $p$ annihilates all ghost-number-$(\neq 0)$ elements.
Since the output of each arity-$k$ tree has ghost number $\leq -1$,
the projection $(p \otimes p)$ kills every summand.

\emph{Step~4 (perturbative stability).}
The BRST perturbation $\delta$ has ghost degree~$+1$, so the
product $h\delta$ preserves ghost number.  For every generator
$\phi \in \cW(\fg)$: $\delta(\phi) = 0$ ($\cW$-generators are
BRST-closed by construction of the DS reduction), so
$(h\delta)(i(\phi)) = 0$, and by induction
$(h\delta)^n(i(\phi)) = 0$ for all $n \geq 1$.  The corrected
inclusion $i_\infty(\phi) = i(\phi)$ is unperturbed, and the
same primitive coproduct results at every perturbative order.
\end{proof}


\subsection{Physical meaning}
\label{subsec:physical-meaning}

The coproduct $\Delta_z$ of a dg-shifted Yangian encodes the
\emph{splitting} of excitations.  Its strictness or nontriviality
is the algebraic fingerprint of the distinction between fundamental
and composite fields.

\begin{center}
\small
\renewcommand{\arraystretch}{1.3}
\begin{tabular}{llccl}
\textbf{Theory}
& \textbf{Boundary algebra}
& $\boldsymbol{\Delta_z}$
& \textbf{Shadow depth}
& \textbf{Reason} \\
\hline
Chern--Simons &
  $\hat{\mathfrak{g}}_k$ (affine KM) &
  $\neq \varepsilon$ &
  $3$ (class $\mathbf{L}$) &
  fundamental gluons split \\[4pt]
3d gravity &
  $\mathrm{Vir}_c = W_k(\mathfrak{sl}_2)$ &
  $= \varepsilon$ &
  $\infty$ (class $\mathbf{M}$) &
  ghost obstruction (DS) \\[4pt]
Twisted holography &
  $H_k$ (Heisenberg) &
  $= \varepsilon$ &
  $2$ (class $\mathbf{G}$) &
  free field, $r = k/z$ \\[4pt]
M2 brane &
  DDCA &
  $\neq \varepsilon$ &
  --- &
  no DS, fundamental modes \\[4pt]
M5 brane &
  $W_{1+\infty}(K)$ &
  $= \varepsilon$ &
  $\infty$ (class $\mathbf{M}$) &
  ghost obstruction (DS) \\
\end{tabular}
\end{center}

The pattern is clean: \emph{theories arising from principal
Drinfeld--Sokolov reduction have primitive coproducts}.  The
ghost sector of the BRST complex provides the grading obstruction
that kills all higher-arity coproduct components.
Theorem~\ref{thm:grav-primitivity-standalone} proves this for
all principal DS reductions $\widehat{\fg} \to \cW(\fg)$: the
ghost-number grading of the BRST complex is universal.
Theories with fundamental excitations (Chern--Simons, deformed
double current algebras) have no ghost sector available, and
their coproducts are nontrivial.

The two exceptions to the DS pattern are illuminating.  The
Heisenberg algebra has a primitive coproduct for a different reason:
it is a free field, its $r$-matrix $r^H(z) = k/z$ is scalar, and
the bar complex is formal at arity~$2$ (shadow depth $r_{\max} = 2$,
class~$\mathbf{G}$).  No DS reduction and no ghost obstruction;
primitivity follows from the absence of nonlinear OPE structure.
The M2 brane boundary algebra (a deformed double current algebra)
has a nontrivial coproduct because it is not obtained by DS
reduction from any affine algebra---its excitations are fundamental
modes of the M2 worldvolume theory.


\subsection{The \texorpdfstring{$r$}{r}-matrix carries all
gravitational structure}
\label{subsec:r-matrix-gravity}

With the coproduct primitive, the full gravitational content of
the Virasoro dg-shifted Yangian is encoded in the $r$-matrix
\begin{equation}\label{eq:grav-r-matrix}
r^{\mathrm{Vir}}(z) = \frac{c/2}{z^3} + \frac{2T}{z}.
\end{equation}
This two-term expression carries a remarkable amount of structure:
\begin{itemize}
\item The \emph{leading coefficient} $c/2 = \kappa(\mathrm{Vir}_c)$
  is the modular characteristic.  It determines the genus tower
  $F_g = (c/2) \int_{\overline{\mathcal{M}}_g} \lambda_g$ at all
  genera, the BTZ black hole entropy
  $S_{\mathrm{BTZ}} = 2\pi\sqrt{ch/6}$, and the Hawking--Page
  transition at $h = c/24$.

\item The \emph{subleading term} $2T/z$ encodes the state dependence
  of gravitational scattering.  The braiding of Wilson lines
  (gravitational line operators) in the holomorphic-topological
  bulk is controlled by this term.

\item The \emph{pole structure}---cubic leading pole, simple
  subleading pole, no even poles---is the signature of a
  bosonic algebra extracted through the $d\!\log$ kernel: the
  OPE poles $z^{-4}, z^{-2}, z^{-1}$ become $r$-matrix poles
  $z^{-3}, z^{-1}$ (each shifted by $-1$), with the $z^{-2}$
  OPE pole producing $z^{-1}$ (which merges with the subleading
  term from $z^{-1}$).
\end{itemize}

The contrast with Chern--Simons is maximal.  In Chern--Simons
gauge theory, the affine $r$-matrix $r^{\mathrm{aff}}(z) = \Omega/z$
is a simple pole and the coproduct is nontrivial: gluons split
and the splitting is nontrivial.  In gravity, the $r$-matrix has
a cubic pole and the coproduct is trivial: gravitons braid but
do not split.  The Drinfeld--Sokolov reduction bridges the two:
it takes the affine Casimir $\Omega/z$ to the gravitational
$(c/2)/z^3 + 2T/z$ by homological perturbation through the
Sugawara nonlinearity, and in the same stroke, the ghost-number
obstruction kills the nontrivial coproduct.

\begin{center}
\small
\renewcommand{\arraystretch}{1.15}
\begin{tabular}{lll}
\textbf{Datum} & \textbf{Chern--Simons} & \textbf{Gravity} \\
\hline
boundary algebra
  & $V_k(\mathfrak{g})$ (affine KM)
  & $\mathrm{Vir}_c$ (Virasoro) \\[2pt]
highest OPE pole
  & $z^{-2}$ (quadratic)
  & $z^{-4}$ (quartic) \\[2pt]
$A_\infty$ tower
  & strict ($m_{k \geq 3} = 0$)
  & infinite ($m_k \neq 0$ all $k$) \\[2pt]
shadow depth
  & $3$ (class $\mathbf{L}$)
  & $\infty$ (class $\mathbf{M}$) \\[2pt]
$r^{\mathrm{coll}}(z)$
  & $\Omega/z$ (simple pole)
  & $(c/2)/z^3 + 2T/z$ (cubic pole) \\[2pt]
$\Delta_z$
  & nontrivial (gluons split)
  & primitive (gravitons braid) \\[2pt]
$\kappa$
  & $\dim\mathfrak{g}\,(k+h^\vee)/(2h^\vee)$
  & $c/2$ \\[2pt]
$\kappa + \kappa^!$
  & $0$
  & $13$
\end{tabular}
\end{center}

The entire table descends from a single datum: the number of
$\lambda$-powers in the $\lambda$-bracket.  Affine Kac--Moody has
$\lambda^1$ at most (quadratic pole), yielding strict associativity,
finite shadow depth, a simple-pole $r$-matrix, and a nontrivial
coproduct.  The Virasoro algebra has $\lambda^3$ (quartic pole),
yielding the full infinite $A_\infty$ tower, infinite shadow depth,
a cubic-pole $r$-matrix, and a primitive coproduct.  The
Drinfeld--Sokolov reduction is the functor that connects the two
columns: it increases the $\lambda$-bracket order by two, trades
the nontrivial coproduct for a ghost obstruction, and replaces the
complementarity $\kappa + \kappa^! = 0$ with $\kappa + \kappa^! = 13$.

\smallskip

In one sentence:
\emph{gravity does not produce particles, only braids; the graviton
is composite, and the ghost sector of its composite origin is the
algebraic obstruction that kills the splitting}.



% ================================================================
%  The E_8 decisive test
% ================================================================

\section{The $E_8$ decisive test}
\label{sec:e8-test}

The bar complex makes a prediction for every lattice vertex
algebra at every genus.  At genus~$2$ for the $E_8$ lattice,
this prediction meets classical number theory.

\subsection{The prediction}

The $E_8$ lattice is even unimodular of rank~$8$.  The vertex
algebra $V_{E_8}$ is class~$\mathbf{G}$ (Gaussian: freely
generated, $\alpha = 0$, $S_4 = 0$) with $\kappa = 8$.
The genus-$2$ formula gives:
\begin{equation}\label{eq:e8-f2}
  F_2(V_{E_8})
  \;=\;
  8 \cdot \frac{7}{5760}
  \;=\;
  \frac{7}{720}.
\end{equation}

\subsection{Verification via Siegel--Weil}

By the Siegel--Weil formula, the genus-$2$ theta function of
any even unimodular lattice of rank~$r$ equals the genus-$2$
Siegel Eisenstein series:
$\Theta_\Lambda^{(2)}(\Omega) = E_{r/2}^{(2)}(\Omega)$.
For $E_8$: $\Theta_{E_8}^{(2)} = E_4^{(2)}$.

The Fourier coefficient $a(T)$ for
$T = I_2 = \bigl(\begin{smallmatrix} 1 & 0 \\
0 & 1 \end{smallmatrix}\bigr)$ counts pairs of orthogonal
roots.  The $E_8$ root system has $240$ roots.  For each
root~$v$, the roots of~$E_8$ partition into:
\begin{enumerate}[(i)]
\item $v$ itself: $1$.
\item $-v$: $1$ (inner product~$-2$).
\item Roots $w$ with $\langle v, w \rangle = \pm 1$: $112$
  (the root system of the stabilizer subalgebra~$E_7$).
\item Orthogonal roots: $240 - 1 - 1 - 112 = 126$.
\end{enumerate}
Therefore $a(I_2) = 240 \times 126 = 30240$.

\begin{theorem}[$E_8$ genus-$2$ three-way agreement]
\label{thm:e8-three-way}
The identity $F_2(V_{E_8}) = 7/720$ is verified by:
\begin{enumerate}[\textup{(}a\textup{)}]
\item the bar-complex formula
  $F_2 = \kappa \cdot \lambda_2^{\mathrm{FP}}$;
\item the Siegel--Weil formula with Cohen--Katsurada
  coefficients;
\item direct lattice enumeration \textup{(}$240 \times 126
  = 30240$\textup{)}.
\end{enumerate}
\end{theorem}

\begin{proof}
The genus-$2$ partition function decomposes as
$Z_2 = \Theta_{E_8}^{(2)} \cdot Z_2^{\mathrm{osc}}$, where
$\Theta_{E_8}^{(2)} = E_4^{(2)}$ by Siegel--Weil and
$Z_2^{\mathrm{osc}}$ is the oscillator contribution.
The Mumford isomorphism gives
$F_g^{\mathrm{osc}} = r \cdot \lambda_g^{\mathrm{FP}}$ with
$r = \operatorname{rank}(\Lambda) = 8$, recovering
$F_2 = 8 \cdot 7/5760 = 7/720$.  The lattice count
$240 \times 126$ is the same number expressed as
$a(I_2)$ of the Siegel modular form~$E_4^{(2)}$.
\end{proof}

The bar complex knows, from purely algebraic data (the OPE
and the modular characteristic~$\kappa = 8$), that each of the
$240$ $E_8$ roots has exactly~$126$ orthogonal partners.  No
root system or lattice enumeration enters the algebraic
prediction; it enters only the verification.


% ================================================================
%  The Epstein zeta function of the shadow metric
% ================================================================

\section{The Epstein zeta function of the shadow metric}
\label{sec:epstein}

\subsection{The shadow metric as a binary quadratic form}

For a class~$\mathbf{M}$ algebra at rational central charge~$c$,
the shadow metric
$Q_L(t) = 4\kappa^2 + 12\kappa\alpha\,t
+ (9\alpha^2 + 2\Delta)\,t^2$
is a quadratic polynomial in~$t$ with rational coefficients,
positive leading coefficient ($4\kappa^2 > 0$), and negative
discriminant ($\operatorname{disc}(Q_L) = -32\kappa^2\Delta < 0$
since $\Delta > 0$ for the standard families at positive~$c$).
Interpreting $Q_L$ as a binary quadratic form
$Q_L(m, n) = 4\kappa^2 m^2 + 12\kappa\alpha\, mn
+ (9\alpha^2 + 2\Delta)\, n^2$
on~$\mathbb{Z}^2$, the form is positive definite.

\begin{theorem}[Epstein zeta of the shadow metric]
\label{thm:shadow-epstein}
Let $\cA$ be a class~$\mathbf{M}$ chirally Koszul algebra
at rational~$c$ with shadow metric~$Q_L$.
\begin{enumerate}[label=\textup{(\roman*)}]
\item The \emph{Epstein zeta function}
\begin{equation}\label{eq:shadow-epstein}
  \varepsilon_{Q_L}(s)
  \;:=\;
  \sideset{}{'}\sum_{(m,n) \in \mathbb{Z}^2}
  Q_L(m,n)^{-s},
  \qquad
  \operatorname{Re}(s) > 1,
\end{equation}
converges absolutely for $\operatorname{Re}(s) > 1$ and
extends to a meromorphic function on~$\mathbb{C}$ with a
simple pole at $s = 1$ of residue
$2\pi/\sqrt{|\operatorname{disc}(Q_L)|}$.

\item The completed function
$\Lambda_{Q_L}(s)
:= (|\operatorname{disc}(Q_L)|/(4\pi^2))^{s/2}\,
\Gamma(s)\,\varepsilon_{Q_L}(s)$
satisfies the functional equation
\begin{equation}\label{eq:shadow-fe}
  \Lambda_{Q_L}(s) = \Lambda_{Q_L}(1 - s).
\end{equation}

\item For classes $\mathbf{G}$ and $\mathbf{L}$
\textup{(}$\Delta = 0$\textup{)}, the form~$Q_L$ degenerates
and the Epstein zeta is not defined as a lattice sum of a
positive definite form.
\end{enumerate}
\end{theorem}

\begin{proof}
Classical~\cite{Epstein1903, Terras73}.  The proof uses
Poisson summation on~$\mathbb{Z}^2$: writing
$Q_L(m,n) = (m,n)\,A\,(m,n)^T$ with the matrix
$A = \bigl(\begin{smallmatrix}
  4\kappa^2 & 6\kappa\alpha \\
  6\kappa\alpha & 9\alpha^2 + 2\Delta
\end{smallmatrix}\bigr)$,
the Mellin transform of the theta function
$\vartheta_{Q_L}(\tau) :=
\sum_{(m,n)} e^{-\pi\tau\,Q_L(m,n)}$
and the standard unfolding argument (splitting the integral
at $\tau = 1$) yield the functional equation.
\end{proof}

\subsection{The shadow field}

\begin{definition}[Shadow field]
\label{def:shadow-field}
For a class~$\mathbf{M}$ algebra at rational~$c$, the
\emph{shadow field} is the imaginary quadratic extension
$K_L := \mathbb{Q}(\sqrt{D_0})$,
where $D_0$ is the squarefree part of
$\operatorname{disc}(Q_L) = -32\kappa^2\Delta$.
\end{definition}

When the class number $h(D_0) = 1$, the Epstein zeta factors
through the Dedekind zeta of~$K_L$:
$\varepsilon_{Q_L}(s) = c_0 \cdot \zeta_{K_L}(s)$
for an explicit constant~$c_0$.  The Dedekind zeta then splits:
$\zeta_{K_L}(s) = \zeta(s) \cdot L(s, \chi_d)$,
where $d$ is the fundamental discriminant and $\chi_d$ the
Kronecker character.  The chain
\[
  \cA
  \;\xrightarrow{\;\text{shadow}\;}
  Q_L
  \;\xrightarrow{\;\operatorname{disc}\;}
  D_0
  \;\xrightarrow{\;\mathbb{Q}(\sqrt{\cdot})\;}
  K_L
  \;\xrightarrow{\;\zeta_K\;}
  L(s, \chi_d)
\]
produces a Dirichlet $L$-function from OPE data.

\subsection{The shadow field table}

The shadow field data for the Virasoro algebra at selected
rational values of~$c$ is recorded in
Table~\ref{tab:shadow-fields-vir}
(Section~\ref{sec:four-classes}).  The squarefree kernel~$D_0$
of $\operatorname{disc}(Q_L) = -320c^2/(5c+22)$ determines
$K_L = \mathbb{Q}(\sqrt{D_0})$.  For the Ising model:
$D_0 = -10$, so $K_L = \mathbb{Q}(\sqrt{-10})$, with
class number $h(-10) = 2$.

The functional equation of the Epstein zeta is a spectral
shadow of Koszul duality: it interchanges $c$ and~$c'$,
and the symmetric point $s = 1/2$ of the completed functional
equation corresponds to the self-dual central charge $c = 13$
where $\Vir_c = \Vir_c^!$.

\begin{remark}[The Koszul--Epstein frontier]
\label{rem:koszul-epstein-frontier}
The deepest open question at the intersection of the algebraic
and arithmetic theories: \emph{does the Koszul--Epstein zeta
function $\varepsilon_{Q_L}(s)$ have all its zeros on a
critical line?}  For class~$\mathbf{M}$ algebras,
$\varepsilon_{Q_L}$ is an Epstein zeta function of a
positive definite binary quadratic form.  The Davenport--Heilbronn
theorem~\cite{DavenportHeilbronn36} shows that Epstein zeta
functions of forms with class number $h > 1$ can have zeros
off the critical line.  The shadow field at the Ising model
has discriminant~$-20$, with class number $h(-20) = 2$; thus
the Ising Epstein zeta may have off-critical zeros.  Whether
such zeros have physical meaning (as resonances, instabilities,
or phase transitions in the shadow obstruction tower) is entirely unknown.
The question connects the algebraic engine (bar complex,
Maurer--Cartan, shadow metric) to the deepest problems in
analytic number theory via a single quadratic form.
\end{remark}


% ================================================================
%  Outlook
% ================================================================

\section{Open problems and outlook}\label{sec:outlook}

\subsection{The descent fan}

The MC element $\Theta_\cA$ admits three independent
projections with distinct targets and distinct symmetries:
\begin{enumerate}[(i)]
\item \emph{Categorical.}
  $\cA \to \barB^{\mathrm{ch}}(\cA) \to \cA^!$.
  The Koszul involution $\cA \leftrightarrow \cA^!$.
\item \emph{Spectral.}
  $\cA \to \{S_r\} \to \sqrt{Q_L} \to \varepsilon_{Q_L}(s)$.
  The functional equation $s \leftrightarrow 1 - s$.
\item \emph{Modular.}
  $\cA \to Z_g(\cA) \to S_\cA(v)$.
  The Galois involution on the $L$-function content.
\end{enumerate}
These three projections share $\Theta_\cA$ as common source
but do not compose into a single functor.  The Koszul,
functional-equation, and Galois involutions are three different
involutions on three different spaces.  For lattice VOAs, all
three projections are compatible (the lattice fan is closed);
for class~$\mathbf{M}$ algebras, the compatibility is open.

\subsection{Multi-generator universality}

Whether $F_g = \kappa \cdot \lambda_g^{\mathrm{FP}}$ holds at
genus~$g \ge 2$ for multi-generator algebras remains open.
The genus-$2$ $\cW_3$ computation
(\S\ref{sec:genus2-w3}) exhibits the cross-channel
correction $(c+120)/(16c)$ that must be cancelled by
CohFT $R$-matrix terms if universality holds.
The decisive evidence would be a direct genus-$2$
multi-channel graph sum with $R$-dressed propagators.

\subsection{The operadic complexity conjecture}

The shadow depth $r_{\max}(\cA)$ is conjectured to equal the
$A_\infty$-depth (the smallest~$n$ such that $m_k = 0$ for
$k > n$ in some $A_\infty$ model) and the $L_\infty$-formality
level (the smallest~$n$ such that $\ell_k = 0$ for $k > n$ on
bar cohomology).  This is proved at arities~$2$, $3$, $4$; the
general statement is open.

\subsection{Pixton's ideal}

Class~$\mathbf{M}$ algebras generate an infinite sequence of
tautological classes on $\overline{\cM}_g$ via the MC tower.
Whether these classes generate Pixton's ideal (the kernel of
the tautological ring surjection) is a deep structural
conjecture connecting the shadow obstruction tower to the geometry of
moduli spaces.


% ================================================================
%  References (placeholder)
% ================================================================

\begin{thebibliography}{99}

\bibitem{BD04}
A.~Beilinson and V.~Drinfeld,
\textit{Chiral Algebras},
AMS Colloquium Publications, vol.~51,
American Mathematical Society, Providence, RI, 2004.

\bibitem{BenjaminChang22}
N.~Benjamin and C.-H.~Chang,
Scalar modular bootstrap and zeros of the Riemann zeta function,
arXiv:2208.02259, 2022.

\bibitem{BGS96}
A.~Beilinson, V.~Ginzburg, and W.~Soergel,
Koszul duality patterns in representation theory,
\textit{J.~Amer.\ Math.\ Soc.}~\textbf{9} (1996), 473--527.

\bibitem{BS88}
A.~Beilinson and V.~Schechtman,
Determinant bundles and Virasoro algebras,
\textit{Comm.\ Math.\ Phys.}~\textbf{118} (1988), 651--701.

\bibitem{CG17}
K.~Costello and O.~Gwilliam,
\textit{Factorization Algebras in Quantum Field Theory},
vols.~1--2, Cambridge University Press, 2017/2021.

\bibitem{DavenportHeilbronn36}
H.~Davenport and H.~Heilbronn,
On the zeros of certain Dirichlet series,
\textit{J.~London Math.\ Soc.}~\textbf{11} (1936), 181--185,
307--312.

\bibitem{DS85}
V.~Drinfeld and V.~Sokolov,
Lie algebras and equations of Korteweg--de Vries type,
\textit{J.~Soviet Math.}~\textbf{30} (1985), 1975--2036.

\bibitem{Epstein1903}
P.~Epstein,
Zur Theorie allgemeiner Zetafunktionen,
\textit{Math.\ Ann.}~\textbf{56} (1903), 615--644.

\bibitem{FG12}
J.~Francis and D.~Gaitsgory,
Chiral Koszul duality,
\textit{Selecta Math.\ (N.S.)}~\textbf{18} (2012), 27--87.

\bibitem{FP00}
C.~Faber and R.~Pandharipande,
Hodge integrals and moduli spaces of curves,
in \textit{Moduli of Curves and Abelian Varieties},
Aspects Math.~E33, Vieweg, 2000, pp.~109--130.

\bibitem{FF90}
B.~Feigin and E.~Frenkel,
Quantization of the Drinfeld--Sokolov reduction,
\textit{Phys.\ Lett.\ B}~\textbf{246} (1990), 75--81.

\bibitem{FlajoletSedgewick}
P.~Flajolet and R.~Sedgewick,
\textit{Analytic Combinatorics},
Cambridge University Press, 2009.

\bibitem{KhanZeng25}
A.~Z.~Khan and K.~Zeng,
Poisson vertex algebras and three-dimensional gauge theory,
arXiv:2502.13227, 2025.

\bibitem{KM94}
M.~Kontsevich and Yu.~Manin,
Gromov--Witten classes, quantum cohomology, and enumerative geometry,
\textit{Comm.\ Math.\ Phys.}~\textbf{164} (1994), 525--562.

\bibitem{Kontsevich03}
M.~Kontsevich,
Deformation quantization of Poisson manifolds,
\textit{Lett.\ Math.\ Phys.}~\textbf{66} (2003), 157--216.

\bibitem{KRW03}
V.~Kac, S.-S.~Roan, and M.~Wakimoto,
Quantum reduction for affine superalgebras,
\textit{Comm.\ Math.\ Phys.}~\textbf{241} (2003), 307--342.

\bibitem{Lorgat}
R.~Lorgat,
\textit{Modular Koszul Duality},
in preparation (monograph, two volumes).

\bibitem{LV12}
J.-L.~Loday and B.~Vallette,
\textit{Algebraic Operads},
Grundlehren der Mathematischen Wissenschaften, vol.~346,
Springer, Berlin, 2012.

\bibitem{Mumford83}
D.~Mumford,
Towards an enumerative geometry of the moduli space of curves,
in \textit{Arithmetic and Geometry, Vol.~II},
Progr.\ Math., vol.~36, Birkh\"auser, 1983, pp.~271--328.

\bibitem{Priddy70}
S.~Priddy,
Koszul resolutions,
\textit{Trans.\ Amer.\ Math.\ Soc.}~\textbf{152} (1970), 39--60.

\bibitem{Terras73}
A.~Terras,
Bessel series expansions of the Epstein zeta function
and the functional equation,
\textit{Trans.\ Amer.\ Math.\ Soc.}~\textbf{183} (1973), 477--486.

\end{thebibliography}

\end{document}
